% Commentary for manuscript pages 64--90. Anchors are verified during assembly.
\begin{dcomment}{exposition}{p64}{Three localizations and one reflective subcategory}
Grothendieck wants the homotopy type of a category over $S$ to retain the way it varies over $S$. The distinction between the three localizations is therefore substantive. Write $\mathcal A=\Cat/S$, and let $W_S$ consist of the morphisms whose underlying functor belongs to $W$. The two finer classes are
\[
 W_S^{\mathrm d}=\{u:X\to Y\mid (s\backslash u)\in W\text{ for every }s\in S\},\qquad
 W_S^{\mathrm g}=\{u:X\to Y\mid (u/s)\in W\text{ for every }s\in S\}.
\]
Here $s\backslash X$ has objects $(x,s\to p(x))$, whereas $X/s$ has objects $(x,p(x)\to s)$. Thus an arrow $s\to t$ induces $t\backslash X\to s\backslash X$ and $X/s\to X/t$. The letters $\mathrm d$ and $\mathrm g$ do not record the variance of the resulting coefficient diagram; the first construction is contravariant, the second covariant. This is the same orientation issue isolated in the earlier commentary at I, manuscript p.~47, but here it determines two different relative homotopy categories.

The relative theorem A hypothesis gives $W_S^{\mathrm d},W_S^{\mathrm g}\subseteq W_S$. Consequently the identity of $\mathcal A$ induces the comparison functors $\rho$ and $\sigma$. These are localizations by the universal property of localization in stages. A priori this says nothing about an embedding in the opposite direction. Such an embedding is supplied by the replacement $X\mapsto C(X)$, together with its natural weak equivalence $\eta_X=i_X:X\to C(X)$ and the fact that $C$ carries a coarse equivalence to either kind of fine equivalence. In particular the composite $\rho i$ is isomorphic to the identity because $\eta$ becomes invertible after coarse localization. This observation establishes a retraction; the following pages establish an adjunction and hence full faithfulness.

For the ordinary homotopy localizer $W_\infty$, the modern interpretations are
\[
 \HOT_{W_\infty}(S)\simeq h\operatorname{Fun}(N(S)^{\mathrm{op}},\mathcal S),
 \qquad
 \HOT^{\mathrm{lc}}_{W_\infty}(S)\simeq h(\mathcal S_{/|N(S)|}),
\]
where $\mathcal S$ is the $\infty$-category of spaces and $h$ denotes its homotopy category. The first expression retains noninvertible transport along the arrows of $S$; the second retains a family parametrized by the homotopy type of its classifying space. Cisinski's comparison of localized categories over a base with locally constant simplicial presheaves makes the latter interpretation precise \citep[\S4.1, especially (4.1.2)]{Cisinski2009LocallyConstant}; straightening supplies the diagram interpretation \citep[Theorem~2.2.1.2]{Lurie2009}.

The subscript $/|N(S)|$ is formed before taking the homotopy category. This distinction already matters for $S=BG$, with $G$ a nontrivial discrete group. The endomorphism set of the point $*\to BG$ in $h(\mathcal S_{/BG})$ is $\pi_0\Omega BG=G$: a morphism includes the homotopy that makes it a morphism over $BG$. Its endomorphism set in the ordinary slice $h\mathcal S/BG$ is a singleton. Coarse localization thus retains relative homotopy information even though its weak equivalences are tested on total categories.
\end{dcomment}

\begin{dcomment}{exposition}{p67}{The two triangles and the role of concatenation}
The purpose of the calculation is to turn the replacement into a universal map from an arbitrary relative homotopy type to a locally constant one. To see exactly what remains to prove, put
\[
 D=\mathcal A[\Sigma_0^{-1}],\qquad E=\mathcal A[\Sigma^{-1}],\qquad
 L:D\longrightarrow E,\quad R:E\longrightarrow D,
\]
where $\Sigma_0\subseteq\Sigma$, $L$ comes from the identity, and $R$ from $C$. Assume $C(\Sigma)\subseteq\Sigma_0$ and $\eta_X\in\Sigma$. The natural transformation $\eta$ descends to a candidate unit $1_D\to RL$, while its inverse in $E$ gives a candidate counit $\epsilon:LR\to1_E$. The first triangle is exactly $\eta_X^{-1}\eta_X=1$. The second is
\[
 C(\eta_X)^{-1}\eta_{C(X)}=1
 \quad\hbox{in }D.
\]
This is why the equality on manuscript p.~68 must be proved in the finer localization. Naturality alone gives
$C(\eta_X)\eta_X=\eta_{C(X)}\eta_X$; cancellation here would require $\eta_X$ to be invertible in $D$, which is precisely what fails for a general $X$.

Concatenation provides the missing cancellation without changing the relative problem. If $C(X)$ consists of pairs $(x,c)$ and $C^2(X)$ of triples $(x,c,c')$, the map
\[
 m_X:C^2(X)\longrightarrow C(X),\qquad(x,c,c')\longmapsto(x,c'c)
\]
is over $S$, because it preserves the final endpoint. The strict left and right unit laws give
$m_XC(\eta_X)=1=m_X\eta_{C(X)}$. Since $C(\eta_X)\in\Sigma_0$, two-out-of-three gives $m_X\in\Sigma_0$. After localization both sections are the inverse of $m_X$, and hence they agree. Associativity of concatenation is not used in this argument. The construction of paths and its preservation of the relevant weak equivalences are earlier inputs; the proof on these pages is complete once those inputs are granted.

For completeness, the resulting adjunction can be written without suppressing either map. A morphism $a:Lx\to y$ in $E$ is sent to $R(a)\eta_x:x\to Ry$. A morphism $b:x\to Ry$ in $D$ is sent to $\epsilon_yL(b):Lx\to y$. Naturality and the two triangle identities show that these operations are mutually inverse and natural in $x,y$. Since $\epsilon$ is an isomorphism, $R$ is fully faithful. This is the ordinary criterion for a reflective localization, expressed with the unit and counit rather than with a chosen presentation by fractions; compare \citet[\S4.3]{Riehl2016}.

One can also identify the objects that have been reflected. An object $x$ is in the essential image of $R$ if and only if $\eta_x$ is an isomorphism. In one direction $\eta_{Ry}$ is inverse to $R\epsilon_y$ by the second triangle. In the other, $\eta_x$ itself exhibits $x\simeq RLx$. The later local-constancy criterion is a geometric description of this elementary categorical condition.
\end{dcomment}

\begin{dcomment}{examples}{p70}{A walking arrow separates the two variances}
Grothendieck's notation becomes concrete over the category $S=[1]=(0\to1)$. Let $P_0$ and $P_1$ be the one-object categories over $S$ supported respectively at $0$ and $1$, and let $T=(S\xrightarrow{1}S)$. We take $W=W_\infty$. Both morphisms $P_i\to T$ are coarse weak equivalences: their classifying spaces are a point and an interval.

The comma categories compute what is lost by this coarse test. For $P_0$ one has
\[
 0\backslash P_0=*,\qquad 1\backslash P_0=\varnothing,
 \qquad P_0/0=*,\qquad P_0/1=*.
\]
For $P_1$ the values are
\[
 0\backslash P_1=*,\qquad1\backslash P_1=*,
 \qquad P_1/0=\varnothing,\qquad P_1/1=*.
\]
Every comma category obtained from $T$ is contractible. Therefore $P_0\to T$ belongs to $W_S^{\mathrm g}$ but not to $W_S^{\mathrm d}$, while $P_1\to T$ belongs to $W_S^{\mathrm d}$ but not to $W_S^{\mathrm g}$. The failure is measured by $\varnothing\to*$, so it survives passage to homotopy types and does not depend on a subtle choice of model. This also supplies the promised example in which $X\to C(X)$ cannot be a fine equivalence: for the right theory, $P_0$ has a varying, rather than locally constant, fibre.

More generally, a contravariant coherent diagram over $[1]$ is an arrow $A_1\to A_0$ of spaces. Its reflection into locally constant diagrams is represented by
\[
 (A_1\longrightarrow A_0)\longrightarrow(A_0\xrightarrow{1}A_0),
\]
with components the given map and $1_{A_0}$. To verify the universal property at the correct level, let $K$ be any space. The derived mapping space from the first diagram to the constant diagram $K$ is the homotopy pullback
\[
 \operatorname{Map}(A_0,K)
 \mathop{\times}^{h}_{\operatorname{Map}(A_1,K)}
 \operatorname{Map}(A_1,K)
 \simeq\operatorname{Map}(A_0,K),
\]
where the second map to $\operatorname{Map}(A_1,K)$ is the identity. Thus its connected components give precisely the desired adjunction on homotopy categories. The left theory has the reversed description: a covariant arrow $A_0\to A_1$ reflects to the constant diagram $A_1$.

This example distinguishes a change of variance from an opposite coefficient category. Both reflections are left adjoints to their respective inclusions. Their different endpoint formulas do not turn either inclusion into a left adjoint to its localization.
\end{dcomment}

\begin{dcomment}{status}{p71}{Duality and the scope of the distinctness assertion}
Grothendieck wants one statement to cover the two orientations of a category over $S$. Taking opposite categories does this, but its variance on functors must be kept explicit. A morphism $u:X\to Y$ over $S$ is carried to $u^{\mathrm{op}}:X^{\mathrm{op}}\to Y^{\mathrm{op}}$ over $S^{\mathrm{op}}$, in the same direction. Moreover
\[
 s\backslash X^{\mathrm{op}}\cong(X/s)^{\mathrm{op}}.
\]
Under the self-duality of $W$ used here, this gives an equivalence
\[
 (\Cat/S)[(W_S^{\mathrm g})^{-1}]
 \simeq(\Cat/S^{\mathrm{op}})[(W_{S^{\mathrm{op}}}^{\mathrm d})^{-1}].
\]
There is no outer opposite on the right. The editorial insertion on this page, which proposes $(\HOT_W(S^{\mathrm{op}}))^{\mathrm{op}}$, conflicts with these explicit localization conventions. It is retained as a reading of the edition, but not adopted mathematically. The two displayed Hom formulas on manuscript p.~72 themselves have the correct common orientation: $\rho\dashv i$ and $\sigma\dashv j$. The later hesitation on p.~84 can be settled by this calculation.

The assertion that the localizations are generally different has genuine content, as the walking-arrow example shows. Footnote~46 strengthens it to a conjecture for every $S$ not equivalent to a discrete category; that strengthening is false. Take $S=BG$ for a nontrivial group $G$. For every $X\to BG$, reindexing $s\backslash X$ along an arrow of $BG$ is an isomorphism of categories. Quillen's theorem B \citep[\S1, Theorem~B]{Quillen1973} therefore identifies its nerve with the homotopy fibre of $NX\to NBG$. A weak equivalence $X\to Y$ over $BG$ induces weak equivalences on these homotopy fibres, so $W_S\subseteq W_S^{\mathrm d}$. The reverse inclusion is the relative theorem A implication already used in constructing $\rho$. Hence $W_S=W_S^{\mathrm d}$; passage to opposites gives $W_S=W_S^{\mathrm g}$ as well. Yet $BG$ is not equivalent to a discrete category when $G\ne1$.

Equivalently, every coherent diagram on a groupoid is locally constant because it carries every invertible arrow to an equivalence. Thus the hypotheses on the base, rather than its being merely non-discrete, control whether the reflection loses information. This correction does not affect the comparison theorem or the adjunction.
\end{dcomment}

\begin{dcomment}{exposition}{p73}{Detecting local constancy without discarding coherence}
Grothendieck now asks which relative homotopy types come from the common locally constant category. The test on manuscript p.~74 concerns the underlying diagram of homotopy types, rather than a classification by such diagrams. For $W=W_\infty$ its precise statement is that a coherent diagram $F:N(S)^{\mathrm{op}}\to\mathcal S$ belongs to the image of the locally constant theory if and only if every transition $F(u)$ is an equivalence. Equivalently, its ordinary underlying diagram $S^{\mathrm{op}}\to h\mathcal S$ sends every arrow to an isomorphism. The equivalence of these two tests uses the definition of equivalence in an $\infty$-category: an arrow is an equivalence exactly when it becomes invertible in the homotopy category.

The reason the test is sufficient is the universal property of groupoid completion. Let $q:N(S)\to B S$ invert all arrows, with $BS$ regarded as a space. Restriction along $q^{\mathrm{op}}$ gives a fully faithful functor
\[
 \operatorname{Fun}((BS)^{\mathrm{op}},\mathcal S)
 \longrightarrow\operatorname{Fun}(N(S)^{\mathrm{op}},\mathcal S),
\]
whose image consists exactly of the functors inverting every arrow. This universal property includes natural transformations and all their coherences. After taking homotopy categories, it proves the asserted essential-image criterion. Cisinski develops the corresponding statement directly for locally constant diagrams and shows that their inclusion has both adjoints, including at the level of derivators \citep[Proposition~3.9 and Theorem~3.12]{Cisinski2009LocallyConstant}.

The strict categorical formulation in terms of $s\backslash X$ is compatible with this account. The split fibration $\Phi X$ has fibre $s\backslash X$ at $s$. If its transition functors are in $W$, it is a locally constant fibre diagram, hence belongs to $\ul{Fib}_0S$. Its image in the fine localization is the original object defined by $X$. Conversely, replacing a locally constant representative by a fine-equivalent one preserves the condition that the transition morphisms are isomorphisms after localization. For the customary homotopy localizer, strong saturation then converts that condition back into membership in $W$. If one drops hypotheses and retains an arbitrary class of arrows, this last conversion must be stated with its saturation; it is not a property of formal localization alone. We do not count this as an additional gap under the full homotopical hypotheses of the scholium.

The factorization of the ordinary underlying diagram through $\Pi S$ records only its one-dimensional monodromy. It does not identify the category of locally constant coherent diagrams with $\operatorname{Fun}((\Pi S)^{\mathrm{op}},\Hot_W)$. In particular, the functor to ordinary diagrams need not be faithful or distinguish isomorphism classes. This is a concrete instance of the earlier commentary's distinction between a coherent diagram and its evaluations (I, manuscript pp.~7 and~11). Here the missing data can be visible in the total homotopy type of the family.
\end{dcomment}

\begin{dcomment}{examples}{p74}{A local system beyond ordinary monodromy}
The criterion just stated allows Grothendieck to recognize a locally constant object without claiming to recover it from its underlying monodromy representation. A familiar bundle makes this distinction explicit. Let $S$ be the poset of nonempty proper faces of a tetrahedron, ordered by inclusion. Its nerve is the barycentric subdivision of the boundary of the tetrahedron, so $|N(S)|$ is a two-sphere. Consider over this sphere the trivial circle bundle
\[
 E_0=S^2\times S^1\longrightarrow S^2
\]
and the Hopf bundle
\[
 E_1=S^3\longrightarrow\mathbb CP^1\cong S^2,
 \qquad (z_0,z_1)\longmapsto[z_0:z_1].
\]
The unit circle acts freely on $S^3\subset\mathbb C^2$ by simultaneous multiplication, and its orbits are exactly the fibres. Over the chart $z_0\ne0$, writing $w=z_1/z_0$, a point of the fibre is uniquely
$\lambda(1,w)/(1+|w|^2)^{1/2}$ with $\lambda\in S^1$. The other projective chart gives a second trivialization. Thus $E_1\to S^2$ is indeed a locally trivial circle bundle, not merely a map with the correct set-theoretic fibres. The classical example and its homotopy-sequence calculation appear in \citet[Examples~4.45 and~4.51]{Hatcher2002}.

The two bundles determine locally constant coherent diagrams with fibre of homotopy type $S^1$. One way to construct such diagrams is to apply the homotopy-fibre functor to paths in the base, retaining its coherent composition; straightening then gives a diagram on the chosen categorical model of the sphere. In both cases the ordinary transport diagram factors through the fundamental groupoid of $S^2$. Since the sphere is simply connected, that groupoid is equivalent to the terminal groupoid. Consequently the ordinary diagrams in $\Hot$ are both isomorphic to the constant diagram at $S^1$.

The coherent local systems are nonetheless different. An equivalence over $S^2$ would induce a weak equivalence on total spaces. But $\pi_1(E_0)\cong\mathbb Z$, whereas $\pi_1(E_1)=\pi_1(S^3)=0$. Therefore no such equivalence exists. The long exact homotopy sequence locates the extra information more precisely:
\[
 \pi_2(S^2)\longrightarrow\pi_1(S^1)\longrightarrow\pi_1(E_i)
 \longrightarrow\pi_1(S^2)=0.
\]
For the Hopf bundle the boundary map is an isomorphism $\mathbb Z\to\mathbb Z$, because the adjacent groups $\pi_2(S^3)$ and $\pi_1(S^3)$ vanish. For the trivial bundle it is zero, because the bundle has a section and its homotopy sequence splits at this place. The distinction is a two-dimensional coherence in the family; it cannot be read from transport around elements of $\pi_1(S^2)$, since there are none.

This explains why the word ``local'' must retain its homotopical meaning. Fibres may all be equivalent and ordinary monodromy may be trivial while the family is nontrivial. Shulman's comparison identifies parametrized spaces with locally constant homotopy sheaves and, over a connected base, with actions of its loop space; the loop space retains precisely the higher information erased by replacing it with its set of components \citep[\S\S5--6 and~8]{Shulman2008}.
\end{dcomment}

\begin{dcomment}{exposition}{p76}{Reading the fifteen categories through replacement functors}
The enlarged diagram is intended to organize choices of representatives, not to produce fifteen distinct homotopy theories. There are two independent decisions. One may represent a relative type by an arbitrary category over $S$, a smooth category, a fibration, or a split fibration. Separately one may require its transport to be locally constant. The upper-level categories impose that additional condition. The cartesian squares express intersections of these requirements. ``Perfect'' here means the categorical smooth-and-proper condition used in the manuscript; it does not refer to perfect complexes or to perfect morphisms in a modern derived-geometric sense.

The replacement $\Phi$ explains the entire left side. For $p:X\to S$, set $F_X(s)=s\backslash X$ and let $\Phi X=\int_{S^{\mathrm{op}}}F_X$ be its Grothendieck construction. An object is a triple $(s,x,a:s\to p(x))$. A morphism to $(t,y,b)$ is a pair $(u:s\to t,v:x\to y)$ satisfying $p(v)a=bu$. Projection onto $s$ makes this a split fibration. There is a canonical functor
\[
 j_X:X\longrightarrow\Phi X,\qquad
 x\longmapsto(p(x),x,1_{p(x)}).
\]
Forgetting $s,a$ defines $r_X:\Phi X\to X$. It satisfies $r_Xj_X=1_X$, and the arrows $(a,1_x)$ give a natural transformation $1_{\Phi X}\to j_Xr_X$. Hence $j_X$ and $r_X$ induce mutually inverse homotopy equivalences of nerves. The retraction $r_X$ need not be over $S$, but $j_X$ is; its underlying weak equivalence proves the coarse comparison. For each $t$, the functor $t\backslash\Phi X\to t\backslash X$ composes the two arrows $t\to s\to p(x)$. Its composite with $t\backslash j_X$ is the identity, and the arrows $(a,1_x)$ again give a natural transformation in the other direction. Thus $t\backslash j_X$ is a homotopy equivalence, proving the fine comparison. In particular, the fibre of $\Phi X$ is exactly the comma category that defines $W_S^{\mathrm d}$.

There is a small but essential issue of morphisms. A split fibration with its chosen splitting corresponds to a strict functor $S^{\mathrm{op}}\to\Cat$. A morphism preserving that splitting corresponds to a strict natural transformation. An arbitrary functor over $S$ between the same total categories need not preserve cartesian arrows, much less the chosen lifts. This is why the four outer arrows in the diagram are not full inclusions. The canonical map $j_X$ cannot by itself be treated as a morphism in every split-fibration category. The comparison theorem uses $\Phi$ and the established replacement maps in the appropriate categories; it does not silently enlarge their morphism sets.

For a smooth $X\to S$, the inclusion $X_s\to s\backslash X$ is a weak equivalence, so the comma replacement has the intended fibrewise value. For a fibration this can be seen directly from cartesian lifting: pulling $(x,a:s\to p(x))$ back along $a$ is right adjoint to the inclusion of $X_s$, and the adjunction supplies the homotopy equivalence. The dual construction $\Psi$ uses $X/s$, producing a split opfibration. Smoothness and properness as intrinsic base-change conditions, and their relationship with categorical fibrations, are developed systematically by \citet[\S\S1 and~3]{Maltsiniotis2005Asphericite}.

The two replacements stay within the upper level when the transitions are weak equivalences. Thus $\Phi_0$ and $\Psi_0$ organize the upper horizontal comparisons. The remaining vertical comparison is the work of $C$, which imposes local constancy and lands in $\ul{Parf}\,S$. The reader can therefore follow the large diagram by tracking three constructions: comma replacement on either side and path replacement between the two levels.
\end{dcomment}

\begin{dcomment}{context}{p81}{Integration and the geometry of a family}
The comparison formula on this page explains why Grothendieck has spent so much effort producing a locally constant representative. If $Y\to S$ is already homotopy-fibred, then the morphisms into $Y$ may be computed in either fine localization or in the coarse localization:
\[
 \Hom_{W_S^{\mathrm d}}(X,Y)\cong
 \Hom_{W_S}(X,Y)\cong
 \Hom_{W_S^{\mathrm g}}(X,Y).
\]
Indeed $Y\to CY$ is a fine equivalence in both senses, and each adjunction identifies maps to $CY$ with maps from the coarse image of $X$ to $Y$. There is no parallel assertion here for maps \emph{out of} an arbitrary locally constant object: the direction matters because these are reflective, rather than merely equivalent, subcategories.

This is part of the project of integrating homotopy types that Grothendieck had described in \emph{Pursuing Stacks}. In its \S69 he connects the inadequacy of a single derived category for trace and determinant constructions with the need to retain derived categories of all diagram shapes \citep{Grothendieck1983}. The present calculation gives this programme a spatial form: a coefficient object varies over an indexing category, and replacing that variation by one along invertible paths is a universal operation. The broad motivation is already discussed in the commentary to I, manuscript p.~4; the new point here is the explicit reflector and its two comma-category presentations.

The word ``coarse'' must also be read relatively. A homotopy equivalence of total spaces over a fixed base is a robust notion only after the base relation has been kept during localization. For a genuine fibration, the long exact sequence of homotopy groups, applied componentwise, implies that a weak equivalence of total spaces over the same base is a weak equivalence on homotopy fibres. This is the geometric model of the axiom L5 invoked in the scholium. It explains why the upper-level categories carry a single weak-equivalence class while the lower-level categories may carry two.

The restriction on $W$ is real. The full chapter's annotation at manuscript p.~43 asserts that the localizer axioms together with L5 force the ordinary homotopy localizer $W_\infty$. Its intended nondegeneracy, expressed nearby through detection of connected components, matters: if one retains only containment $W_\infty\subseteq W$, the maximal class of all functors satisfies the formally stated axioms and L5 as well. The later statement makes this explicit: Chapter~XII, Theorem~1 on manuscript p.~12, assumes the connected-components axiom, and notes~27--28 allow the weaker exclusion $W\ne\mathrm{Fl}(\Cat)$ \citep{Grothendieck1990XIIa}. We therefore use the concrete ordinary homotopy localizer for the modern identification, and keep L5 explicit in abstract applications. In particular one cannot read the conclusion as a theorem for every homological or cohomological localization. For example, a degree-$n$ covering of a circle can be a homology equivalence with coefficients in a field in which $n$ is invertible, while its fibres over the target circle have different numbers of components from those of the identity. The fibrewise inference underlying L5 then fails. The theorem's conditional formulation in terms of L5 is correct; its apparent breadth in the letter $W$ should not conceal that condition.
\end{dcomment}

\begin{dcomment}{later}{p83}{What extra geometry a relative theory needs}
Grothendieck is separating two questions that a model structure alone can easily conflate. A model category supplies a homotopy theory of objects over $S$, with weak equivalences tested on their total objects. To obtain the finer theory contemplated here, one must also say what the local probes of $S$ are and how those probes retain directed variation. For a small category, the two comma constructions supply such probes and distinguish the two variances.

Thomason's model structure solves the first problem for $\Cat$: its weak equivalences are nerve weak equivalences, and its fibrations are detected by $\operatorname{Ex}^2N$; see Theorem~4.9 and Corollary~5.5 of \citet{Thomason1980}. This gives powerful factorization and invariance results, but it does not by itself identify a Thomason fibration with one of Grothendieck's smooth or homotopy-fibred representatives. Nor does it manufacture a second directed localizer from the bare triple of classes. The categorical comma constructions remain additional geometry.

In the later theory, left and right fibrations over an $\infty$-category provide the two space-valued diagram theories, and straightening identifies them with covariant and contravariant functors. When the base is already an $\infty$-groupoid, every arrow is invertible and these theories describe locally constant families. This explains the intuition behind Grothendieck's remark about $\infty$-groupoids, without claiming that they admit no notion of relative homotopy or localization. Their slices are rich; what is absent is the noninvertible transport of a general indexing category.

The comparisons between the named ``Zariski'', smooth, proper, and dual topologies are not formal consequences of these names. Precise statements require specifying covering families and the associated class of equivalences. Maltsiniotis's asphericity structures supply a rigorous framework for the categorical smoothness and base-change part of this discussion. We therefore treat the paragraph as a programme of comparison, not as an unproved theorem asserting that all four resulting cohomology theories agree.

The distinction is obscured by a paired reversal of superscripts in the long
paragraph on manuscript p.~82. As the definitions and the explicit statement
on p.~94 show, local-isomorphism tests describe $W_S^{\mathrm g}$, whereas
co-local-isomorphism tests describe $W_S^{\mathrm d}$. The edition prints the
opposite association in that paragraph. We preserve its wording and use the
definitions throughout.

\end{dcomment}

\begin{dcomment}{development}{p84}{The coherence also holds after universal localization}
Grothendieck asks whether the path-replacement coherence remains true for the finer class of equivalences stable under every change of base. His homotopy-fibred representatives allow an affirmative answer under the hypotheses already used in the scholium. It is useful to isolate the exact argument, since it identifies where L5 enters and gives a sharper intermediate localization.

\textbf{Universal replacement proposition.}\ours\ Let $W$ be a class of functors containing all isomorphisms and satisfying two-out-of-three, and put
\[
 U_S=\{u:X\to Y\text{ over }S\mid
 T\times_Su\in W\text{ for every }T\to S\}.
\]
For every base $T$, suppose a full subcategory $\mathcal B_T\subseteq\Cat/T$ is given, preserved by arbitrary base change, such that for morphisms between its objects membership in $W_T$ is equivalent to membership in $W_T^{\mathrm f}$. Suppose $C:\Cat/S\to\mathcal B_S$ has a natural map $\eta:1\to C$ in $W_S$, and that $C$ preserves $W_S$. Finally suppose $m:C^2\to C$ satisfies
\[
 m_X\eta_{C(X)}=1=m_XC(\eta_X).
\]
Then $C(W_S)\subseteq U_S$; the equality $\eta_{C(X)}=C(\eta_X)$ holds in $(\Cat/S)[U_S^{-1}]$; and localization induces an adjunction
\[
 (\Cat/S)[U_S^{-1}]
 \mathrel{\mathop{\rightleftarrows}^{L}_{R}}
 (\Cat/S)[W_S^{-1}],\qquad L\dashv R,
\]
where $R$ is induced by $C$ and is fully faithful.

\emph{Proof.} Let $v:A\to B$ be a $W_S$-morphism between objects of $\mathcal B_S$. By hypothesis each fibre $v_s$ belongs to $W$. For $a:T\to S$, the objects $a^*A,a^*B$ belong to $\mathcal B_T$, and the fibre of $a^*v$ at $t$ is exactly $v_{a(t)}$. The converse fibre criterion in $\mathcal B_T$ gives $a^*v\in W$. Since $a$ was arbitrary, $v\in U_S$. Apply this to $C(u)$ for $u\in W_S$: its source and target lie in $\mathcal B_S$, and it lies in $W_S$ by assumption. This proves $C(W_S)\subseteq U_S$.

The class $U_S$ contains the identities and satisfies two-out-of-three, since each pullback functor does and $W$ does. It is contained in $W_S$ by taking $T=S$. In particular $C(\eta_X)\in U_S$. From $m_XC(\eta_X)=1$, two-out-of-three yields $m_X\in U_S$. The two unit sections of $m_X$ therefore become its inverse in the $U_S$-localization, proving their equality. The unit and counit constructed in the preceding adjunction calculation now apply with $\Sigma_0=U_S$ and $\Sigma=W_S$; their triangles were proved explicitly there. The counit is the inverse of the coarse image of $\eta$, and hence is invertible. This proves the remaining assertions.\hfill$\square$

In Grothendieck's application, take $\mathcal B_T=\ul{Fibhot}\,T$ and $m$ to be path concatenation. Arbitrary base-change stability and the agreement of fibrewise and coarse equivalence are exactly the hypotheses used above; the latter includes L5. Also $U_S\subseteq W_S^{\mathrm d}\cap W_S^{\mathrm g}$, since the comma tests are the pullbacks along $s\backslash S\to S$ and $S/s\to S$. Thus the result strengthens the specific coherence calculation, without asserting that either fine localization equals the universal one.

We have checked the general comparison machinery in Thomason and the treatments of locally constant diagrams in Cisinski and Shulman cited above. We could not find this exact intermediate reflection stated there in this form. It is an extraction from the manuscript's hypotheses, not a priority claim; the more general reflection argument attached below to the second scholium explains why concatenation can ultimately be dispensed with.
\end{dcomment}

\begin{dcomment}{exposition}{p87}{Why replacements of objects must also control morphisms}
Grothendieck now removes the concrete path construction to find what the comparison actually uses. Choosing a weak equivalence $X\to C(X)$ with fibrant target proves that the inclusion of fibrant objects is essentially surjective after localization. It does not by itself identify morphisms between those objects. This is the source of the two alternatives he introduces: a functorial replacement, or a lifting argument that produces a functor only after passing to homotopy.

The functorial alternative is a complete and economical proof. Let $D:\mathcal B\hookrightarrow\mathcal A$ be full, let $C:\mathcal A\to\mathcal B$, and let $\eta:1_{\mathcal A}\to DC$ be natural with components in $W$. Suppose $W$ satisfies two-out-of-three and give $\mathcal B$ the restricted class. For $u:X\to Y$ in $W$, naturality says
\[
 DC(u)\eta_X=\eta_Yu.
\]
The right-hand side and $\eta_X$ lie in $W$, so $DC(u)\in W$. Hence $C$ descends to the localizations. The maps $\eta_X$ give $1_{\mathcal A[W^{-1}]}\cong\bar D\bar C$. Restricting $\eta$ to the full subcategory $\mathcal B$ gives $1_{\mathcal B[W^{-1}]}\cong\bar C\bar D$. Thus the two localized functors are quasi-inverse equivalences. No triangle identity is needed to establish an equivalence from these two natural isomorphisms.

If $\mathcal B\subseteq\mathcal F\subseteq\mathcal A$ with $\mathcal F$ full, the same proof restricts $C$ and $\eta$ to $\mathcal F$. Both inclusions therefore induce equivalences. This proves the intermediate-subcategory assertion on p.~89, including full faithfulness, rather than only the existence of representatives. It also makes the independence of choices precise. Two replacement systems determine a unique natural comparison \emph{compatible with their chosen comparison isomorphisms to the identity}. Without this compatibility clause, quasi-inverse functors may have nontrivial automorphisms, so ``unique up to unique isomorphism'' is not literally a statement about bare functors.

In the lifting alternative, each $i_X$ must have the \emph{left} lifting property against the fibrations. The phrase ``relèvement à droite'' on p.~86 is a slip of the edition, confirmed by the square on p.~90: $i_X$ is its left vertical arrow and $C(X')\to S$ is its right vertical fibration. A filler $\bar u:C(X)\to C(X')$ extends $i_{X'}u$. The existence of the filler gives a candidate image of $u$; independence of the filler, identity preservation, and composition are the work of the lemma on the following pages. That argument is logically distinct from the elementary functorial proof above. Keeping this distinction visible explains why the second scholium has two routes to the same comparison theorem.
\end{dcomment}

\begin{dcomment}{exposition}{p91}{A path object supplies uniqueness of the lift}
Grothendieck wants to replace every object over $S$ by one in a useful full subcategory, without making a functorial choice of the replacements themselves. The point of the next argument is that functoriality can appear after localization even when it is absent before localization. A lift need not be unique as a map; its class must be unique in precisely the category in which the replacement is to live. This distinction is the work performed by the relative path object on pp.~92--93.

Write $\mathcal A=\mathcal M/S$ and $\mathcal F=\underline{\mathrm{Fib}}\,S$. The class $\mathrm{Cof}_0$ consists of maps having the \emph{left} lifting property against every member of $\mathrm{Fib}$. In this generality the notation does not assert $\mathrm{Cof}_0\subseteq W$; the argument assumes separately that the particular map under consideration belongs to $W$. Suppose that $i:X\to C$ belongs to $\mathrm{Cof}_0\cap W$, that $C,D$ lie in $\mathcal F$, and that two maps $v_0,v_1:C\to D$ over $S$ agree after composition with $i$. Factor the relative diagonal as
\[
D\xrightarrow{j}P_SD\xrightarrow{(d_0,d_1)}D\times_S D,
\qquad j\in W,\quad (d_0,d_1)\in\mathrm{Fib}.
\]
The target of the diagonal is the fibre product over $S$, as the edition already corrects. The two maps to $D$ agree on $X$, so the lifting property supplies $h:C\to P_SD$ with $d_0h=v_0$ and $d_1h=v_1$. This is a right homotopy relative to $X$.

The argument must remain inside $\mathcal F$. Since $D\to S$ is a fibration, the projection $D\times_S D\to D$ is a fibration by base change; composition then shows $D\times_S D\to S$ is a fibration. Composing once more with $(d_0,d_1)$ proves that $P_SD$ belongs to $\mathcal F$. Thus $j$ can be inverted in $\mathcal F[W^{-1}]$. The equalities $d_0j=1_D=d_1j$ imply $[d_0]=[d_1]$, whence $[v_0]=[v_1]$. No assertion about a calculus of fractions, and no cofibrancy of $C$, is required. In fact this cancellation proof does not use $i\in W$: the lifting property alone supplies $h$. Membership in $W$ is needed later for the maps $i_X$ to give replacement isomorphisms. The path object is used only to transport the equality into the correct localization.

This supplies all the functoriality left implicit in the text. Choose for each $X$ a factorization $X\xrightarrow{i_X}C_X\to S$ with $i_X\in\mathrm{Cof}_0\cap W$ and $C_X\in\mathcal F$. For $u:X\to Y$, choose a lift $\bar u:C_X\to C_Y$ satisfying
\[
\bar u i_X=i_Yu.
\]
Such a lift exists by applying the lifting property to $C_Y\to S$, and the preceding argument makes $[\bar u]$ independent of its choice. The identity of $C_X$ is an allowable lift for $1_X$. If $u:X\to Y$ and $v:Y\to Z$, the composite $\bar v\bar u$ is an allowable lift of $vu$. Consequently $\gamma(X)=C_X$, $\gamma(u)=[\bar u]$ defines a functor to $\mathcal F[W^{-1}]$. When $u\in W$, the equation above and two-out-of-three give $\bar u\in W$; hence $\gamma$ descends to $\mathcal A[W^{-1}]$.

Finally, the same equation proves naturality of $i_X$, first for actual arrows and then for formal inverses. For an intermediate full subcategory $\mathcal F\subseteq\mathcal E\subseteq\mathcal A$, the maps $i_X$ therefore give natural isomorphisms between both composites of the inclusion and $\gamma$ and the relevant identity functors. This proves the comparison equivalences claimed on p.~94, including full faithfulness. The unproved words ``\`A v\'erifier la fonctorialit\'e'' ask for this short verification, rather than for an additional homotopical theorem. The same method reappears in the replacement theorem for derivable categories: \citet[Proposition~1.8]{Cisinski2010} constructs the quasi-inverse under a different, weaker organization of the axioms.
\end{dcomment}

\begin{dcomment}{exposition}{p94}{What is being localized relatively}
Grothendieck now separates information about the total object from information that remains visible after testing over smaller bases. The distinction is needed because a weak equivalence of total categories can conceal a change in its fibres. Given the test maps $t:T\to S$ belonging to $\mathcal L$, the definition is simply
\[
W_S^{\mathcal L}=\bigcap_{t\in\mathcal L/S}(t^*)^{-1}(W_T).
\]
Each condition concerns an ordinary pullback before any localization. Consequently $W_S^{\mathcal L}$ contains the isomorphisms and satisfies two-out-of-three whenever $W$ does. If $1_S$ is a test map, it is contained in $W_S$. Neither closure of $\mathcal L$ under composition nor closure under base change is needed for these assertions. Grothendieck's hesitation in footnote~57 can therefore be answered affirmatively for the construction at a fixed base: composition closure is dispensable there.

Composition becomes relevant when bases vary. If $g:T\to S$ belongs to $\mathcal L$ and $\mathcal L$ is composition closed, then $g^*$ carries $W_S^{\mathcal L}$ into $W_T^{\mathcal L}$. Indeed, for every test map $h:U\to T$, the composite $gh$ is a test map over $S$, and $h^*g^*f$ identifies with $(gh)^*f$. If $g$ is arbitrary, the same proof requires the more specific condition that every such composite $gh$ be an allowed test over $S$. A theory at one base and a contravariant theory of all bases thus require different amounts of structure.

The inclusion of weak-equivalence classes gives the canonical localization
\[
\rho:\mathcal A[(W_S^{\mathcal L})^{-1}]\longrightarrow
\mathcal A[W_S^{-1}].
\]
This assertion does not require either class to be saturated: the universal property says that successively inverting the images of $W_S$ has the same effect as inverting $W_S$ initially. If ``localization functor'' is defined using the entire inverse image of the isomorphisms, rather than the nominated class $W_S$, the conclusion is unchanged, since saturation does not change the localized category.

Hypothesis~1 is the bridge between these two localizations. It says that the fine and coarse classes have identical restrictions to the full subcategory of fibrations over $S$. The replacement equivalence then supplies a canonical functor, up to the usual unique isomorphism,
\[
i:\mathcal A[W_S^{-1}]\longrightarrow
\mathcal A[(W_S^{\mathcal L})^{-1}],
\qquad X\longmapsto C_X,
\]
with $\rho i\simeq1$. The latter equation alone does not make $i$ fully faithful: a functor can have a retraction without being full. What turns the retraction into an adjunction is the natural map $X\to C_X$ in the \emph{fine} localization and its compatibility on replaced objects. This is why the triangle identity returns on pp.~98--100. The terminology ``locally constant'' for the coarse localization has geometric content in the categorical application; the abstract axioms, by themselves, provide a category of fibrant objects over $S$, without a separately defined notion of local constancy.
\end{dcomment}

\begin{dcomment}{exposition}{p97}{Base change between fibrant objects}
The lemma on pp.~95--97 protects the replacement subcategory from the loss of information caused by strict pullback. Its hypotheses concern weak equivalences between two fibrations over the same base. They do not assert that pullback preserves every weak equivalence over that base.

Here is the proof with the two changes of base kept separate. Let $w:X\to Y$ over $S$ lie in $W$, with $X\to S$ and $Y\to S$ fibrations, and factor an arbitrary $g:T\to S$ as $T\xrightarrow{a}\bar T\xrightarrow{b}S$, where $a\in W$ and $b\in\mathrm{Fib}$. Pullback along $b$ preserves $w$: the map $Y\times_S\bar T\to Y$ is a fibration, and right properness applies to the square obtained by pulling $w$ back along that map. Both new structure maps to $\bar T$ are fibrations. Pullback along $a$ makes the two vertical comparison maps
\[
X\times_ST\longrightarrow X\times_S\bar T,
\qquad
Y\times_ST\longrightarrow Y\times_S\bar T
\]
weak equivalences, again by right properness. The commutative square and two-out-of-three then make $g^*w$ a weak equivalence. This proves Hypothesis~1 even when every map is admitted as a test map. The printed condition $g\in\mathrm{Fib}$ in the factorization must concern $b=\bar g$; the edition supplies that correction.

In a full model category the conclusion has an additional explanation, and does not require right properness. Pullback $g^*:\mathcal M/S\to\mathcal M/T$ is right Quillen for the slice model structures, since it preserves fibrations and trivial fibrations. Ken Brown's lemma therefore preserves weak equivalences between fibrant objects of the slice; those objects are exactly the fibrations to $S$ \citep[Lemma~11.3.14 and its dual]{Riehl2014}. One can see the relevant mechanism directly. Factor $(1_X,w):X\to X\times_SY$ as a trivial cofibration $j:X\to Z$ followed by a fibration $(p,q)$. Both $p$ and $q$ are fibrations, because $X$ and $Y$ are fibrant in the slice. The equations $pj=1_X$ and $qj=w$ make both trivial fibrations. Their pullbacks are therefore weak equivalences; $g^*p\,g^*j=1$ first implies that $g^*j$ is a weak equivalence, and $g^*w=g^*q\,g^*j$ proves the assertion. The last calculation only needs $j\in W$, not its cofibration property. Thus it also works under Grothendieck's factorization and fibration-stability assumptions as soon as $\mathrm{Fib}\cap W$ is pullback stable. Hypothesis~4 below already implies Hypothesis~1 for every class of test maps. This is the usual Brown factorization argument applied in a slice. Grothendieck's right-properness hypothesis is another valid sufficient condition, stronger than what this calculation requires.

A small example displays the distinction. In simplicial sets over $\Delta^1$, the map from its initial vertex $\{0\}\to\Delta^1$ to the terminal object $1_{\Delta^1}$ is a weak equivalence of total simplicial sets: both realizations are contractible. Pull it back along the terminal vertex $\{1\}\to\Delta^1$. The resulting map is $\varnothing\to\Delta^0$, which fails even to induce a bijection of connected components. The structure map $\{0\}\to\Delta^1$ is not a Kan fibration, so the lemma does not apply. Replacing it by a fibration before pulling back is precisely what removes the counterexample. In particular, the assertion on p.~95 should never be shortened to ``weak equivalences are stable under arbitrary base change.'' Its restriction to fibrant objects is the mathematical content.
\end{dcomment}

\begin{dcomment}{exposition}{p99}{The unit must live in the finer localization}
Grothendieck is trying to identify a universal map from an arbitrary relative object to a locally constant one. The intended adjunction is $\rho\dashv i$, where $\rho$ forgets the finer relative information and $i$ puts the resulting object back by replacement. Its unit is therefore the map $\eta_X:X\to C_X$ in the finer category. Its counit is the inverse of the same map after passing to the coarser category. This accounts for the direction of the two transformations without a convention about ``left'' and ``right'' fibrations.

In the lifting construction the map $\eta_X$ is natural already in the fine localization. To check this, choose a lift $\bar u$ of $u:X\to Y$. The equality $\bar u\eta_X=\eta_Yu$ holds before localization. A different lift has the same class in $\mathcal F[W^{-1}]$, and that localization maps to $\mathcal A[(W_S^{\mathcal L})^{-1}]$ because the two classes agree on $\mathcal F$. Thus the displayed equation gives a well-defined naturality square in the fine category. Once it holds for actual arrows it holds for inverted weak equivalences as well. The reference to the coarse localization in the prose on p.~98 understates what this construction establishes; the finer assertion needed for the adjunction is available.

The first triangle identity is immediate: after applying $\rho$, the composite is $[\eta_X]^{-1}[\eta_X]$. The second reduces to
\[
[C(\eta_X)]=[\eta_{C_X}]
\quad\hbox{in }\mathcal F[W^{-1}].
\]
The two maps agree after composition with $\eta_X$, and $\eta_X$ has the left lifting property used in the cancellation lemma. Hence the proof on p.~100 establishes this equality, and its image in the fine localization establishes the second triangle. Since the counit is invertible, $i$ is fully faithful. The adjunction gives, functorially in both variables,
\[
\Hom_{\mathcal A[(W_S^{\mathcal L})^{-1}]}(X,C_Y)
\simeq
\Hom_{\mathcal A[W_S^{-1}]}(X,Y).
\]
For a fibrant $Y$, the map $Y\to C_Y$ belongs to both weak-equivalence classes, so the right-hand side is equally the Hom set into $Y$ in the fine localization. This is the corollary eventually stated on p.~113. It compares \emph{localized} Hom sets; it does not identify them with the set of strict maps $X\to Y$.

The distinction between strict choices and their localized classes explains why the proof remains valid without functorial factorizations. It also isolates what the next path-object calculation is supposed to achieve when a strict functor $C$ is available. Below, after Theorem-Scholie~2, we give a formal argument showing that a functorial replacement into a full subcategory already contains enough information for that calculation.
\end{dcomment}

\begin{dcomment}{exposition}{p103}{The mapping-path construction and its hypotheses}
Grothendieck now wants to make the replacement geometric. An object $X\to S$ is enlarged by adjoining paths in the base which start at the image of a point of $X$; the enlarged object is then placed over the other endpoint. This is the reason for keeping both endpoint maps of $\mathrm{Ch}(S)$ in the large diagram. The projection back to $X$ preserves the starting point and need not be a morphism over $S$ when $C_S(X)$ is placed over the endpoint.

Put $P(S)=\mathrm{Ch}(S)$ and write its endpoint maps as $\sigma,\beta:P(S)\rightrightarrows S$. The identity paths give $\alpha:S\to P(S)$, with $\sigma\alpha=\beta\alpha=1$. The construction is
\[
C_S(X)=X\times_{S,\sigma}P(S),\qquad
C_S(X)\longrightarrow S:\ (x,c)\longmapsto\beta(c),
\]
where the notation for elements is shorthand for a statement about all generalized elements. The map $\eta_X$ sends $x$ to $(x,\alpha(f_Xx))$. The two required properties, $C_S(X)\to S\in\mathrm{Fib}$ and $\eta_X\in W$, have different proofs and need different hypotheses.

First, $C_S(X)\to X\times S$ is the pullback of $(\sigma,\beta):P(S)\to S\times S$. It is a fibration by Hypothesis~2 and pullback stability. Its composite with $X\times S\to S$ is a fibration when $X\to e$ is a fibration. Hypothesis~3 imposes this for every $X$, so the construction works for every object of every slice. This is an all-objects-fibrant framework, not a claim that every object of an arbitrary model category is fibrant.

Second, two-out-of-three applied to $\sigma\alpha=1$ and $\beta\alpha=1$ gives $\sigma,\beta\in W$. Both endpoint maps are fibrations, using Hypothesis~3 and closure of fibrations under composition. Hypothesis~4 then says that their pullbacks remain weak equivalences. In particular, the projection $r_X:C_S(X)\to X$, obtained by pulling back $\sigma$, belongs to $W$. Since $r_X\eta_X=1_X$, two-out-of-three gives $\eta_X\in W$. One must use the projection $r_X$ here: directly pulling back $\alpha\in W$ would require a stronger assumption, which has not been made. The two editorial repairs on p.~103, replacing $\mathrm{pr}_1\pi'$ by $\mathrm{pr}_1p'$ and ``dans $i_X$'' by ``dans $W$'', restore precisely this argument.

The language of generalized elements introduces no additional geometry. Yoneda and preservation of finite limits by $\Hom(T,-)$ identify $C_S(X)(T)$ with pairs $(x,c)$ satisfying $\sigma(c)=f_X(x)$, and identify $C_S^2(X)(T)$ with triples $(x,c,c')$ satisfying $\sigma(c')=\beta(c)$. These identifications are natural in $T$, so a formula on all such triples defines a morphism in $\mathcal M$ if it is natural in $T$. The two unit maps are
\[
\begin{aligned}
\eta_{C_S(X)}(x,c)&=(x,c,1_{\beta(c)}),\\
C_S(\eta_X)(x,c)&=(x,1_{f_X(x)},c).
\end{aligned}
\]
The edition gives these formulas in the opposite order after the word ``respectivement'' on p.~105. This is an edition-level transposition, not a failure of the proof: the argument treats the two maps symmetrically. Likewise the $\mathrm{Ch}(T)$ printed on p.~106 must be read $\mathrm{Ch}(X)(T)$ when describing the arrows with object set $X(T)$.

A composition law on paths preserving endpoints gives $\mu_X(x,c,c')=(x,c'c)$. Bilateral identity laws make $\mu_X\eta_{C_S(X)}=1=\mu_XC_S(\eta_X)$. Since $\eta_{C_S(X)}\in W$, two-out-of-three implies $\mu_X\in W$; both its source and target are fibrant over $S$, so it becomes invertible in their localization and equalizes the two unit maps there. Associativity of composition is indeed unused. A strictly associative path construction contains more structure than this particular adjunction requires.
\end{dcomment}

\begin{dcomment}{examples}{p107}{The whole construction over a groupoid}
The geometric meaning of the calculation can be seen without topological path spaces. Take $\mathcal M$ to be the category of small groupoids, let $W$ be the equivalences of categories, and let $\mathrm{Fib}$ be the isofibrations. An isofibration lifts any isomorphism in the base with prescribed source. This setting gives a complete example of the path construction, including the distinction between an equivalence of total objects and an equivalence after base change. In this example the fine class is $U_S=W_S^{\mathrm{Fl}(\mathcal M)}$, tested by \emph{all strict changes of base}. It must be kept distinct from the directional class $W_S^{\mathrm d}$ in $\mathrm{Cat}/S$, whose testing maps form a different family; the calculation below does not identify these two classes.

For a groupoid $S$, let $P(S)$ be its arrow groupoid. Its objects are arrows $c:s\to t$ of $S$; an arrow from $c$ to $c':s'\to t'$ is a commuting square, specified by isomorphisms $a:s\to s'$ and $b:t\to t'$ with $bc=c'a$. The endpoint functor sends the square to $(a,b)$. It is an isofibration: given $c$ and any pair $(a,b)$, the arrow $c'=bca^{-1}$ supplies the desired lift. The identity-arrow functor $\alpha:S\to P(S)$ is fully faithful, since a square between identity arrows is determined by its common vertical arrow. It is essentially surjective because every $c:s\to t$ is isomorphic to $1_s$, by the square $(1_s,c)$. Thus $\alpha\in W$. Every groupoid maps by an isofibration to the terminal groupoid, so Hypotheses~2 and~3 hold.

For completeness, a trivial fibration $u:E\to B$ of groupoids stays a trivial fibration after arbitrary pullback. It stays an isofibration by the lifting definition. An equivalence of groupoids which is an isofibration is surjective on objects: essential surjectivity supplies an isomorphism to an image object, and the isofibration property lifts that isomorphism to an object with exactly the prescribed image. A pullback of $u$ is consequently surjective on objects and fully faithful; the latter follows by applying full faithfulness to the equality of the prescribed images of an arrow. This proves Hypothesis~4 directly.

Now let $f:X\to S$ be any functor between groupoids. The groupoid $C_S(X)$ has objects $(x,c)$ with $c:f(x)\to s$. An arrow $(x,c)\to(y,d)$ is a pair $(u,v)$ satisfying
\[
u:x\to y,\qquad v:s\to t,\qquad vc=d\,f(u).
\]
Its structure functor to $S$ sends $(x,c)$ to $s$ and $(u,v)$ to $v$. It is an isofibration, since an arrow $v:s\to t$ lifts from $(x,c)$ to $(x,vc)$ by $(1_x,v)$. The projection $r:C_S(X)\to X$ has section $\eta(x)=(x,1_{f(x)})$. There is a natural isomorphism $\eta r\Rightarrow1_{C_S(X)}$ whose component at $(x,c)$ is $(1_x,c)$. Hence $r$ and $\eta$ are equivalences of total groupoids. Notice that $r$ is generally not over $S$: it sends $(x,c)$ to an object over $f(x)$, whereas the structure map of $C_S(X)$ sends it to the target of $c$.

An object of $C_S^2(X)$ is a triple $(x,c,c')$. The multiplication is
\[
\mu(x,c,c')=(x,c'c).
\]
If a morphism of triples is given by $(u,v,w)$ with $vc=df(u)$ and $wc'=d'v$, its image is $(u,w)$; the equality $wc'c=d'df(u)$ verifies functoriality. The two unit equalities hold on both objects and arrows. Since $\mu\eta_C=1$ and $\eta_C$ is an equivalence, $\mu$ is an equivalence. This verifies every ingredient of the localized triangle identity in an explicit categorical model.

Take now $S=BG$ for a nontrivial group $G$, regarded as a one-object groupoid, and take $X=e\to BG$. The objects of $C_{BG}(e)$ are the elements $g\in G$. There is exactly one arrow from $g$ to $h$, labelled in the base by $hg^{-1}$. Thus $C_{BG}(e)$ is the connected groupoid with exactly one arrow between any two objects and is equivalent to $e$. Its projection to $BG$ is nevertheless far from having a one-point strict fibre. Pulling it back along $e\to BG$ retains only the arrows labelled $1$, and produces the discrete groupoid on the set $G$. The pullback of $X$ is $e$, and the pullback of $\eta_X$ is the inclusion of $1\in G$ into this discrete groupoid. For $G\ne1$ it is not an equivalence. The replacement unit is therefore in $W_S$ but not in the class of equivalences stable under all changes of base.

This computation also identifies the fibre one expects homotopically: replacing the point over $BG$ by a fibration reveals the loop group $G$. It explains why the coarser localization may invert a replacement map which the finer one must retain as a noninvertible unit. The source object $e\to BG$ is not an isofibration when $G\ne1$, so the example is consistent with the base-change result for maps between fibrations. If $\mathcal L$ consists of all base maps, the fine and coarse equivalences do agree on isofibrations, by the slice-model argument above or by the same lifting calculation on groupoids. No assertion of literal local triviality of a topological bundle is needed for this example.
\end{dcomment}

\begin{dcomment}{later}{p108}{The axioms among Brown and Quillen}
Grothendieck's abstraction is useful because the localization comparison asks for much less structure than a full model category. In version~B, only one sort of factorization and its associated lifting property enter the construction of replacement; the path-object proof then supplies the equality needed for functoriality. The axioms as stated do not assert a second weak factorization system, closure of weak equivalences under retracts, arbitrary small colimits, or small Hom sets after localization. These should not be silently imported from the familiar model-category terminology.

There is a close earlier precedent in Brown's categories of fibrant objects. Their basic data are weak equivalences, fibrations, pullbacks along fibrations, and path objects, with every object fibrant; Brown develops this in Part~I, \S1 of \citet{Brown1973}. Grothendieck's Hypotheses~2--4 place his mapping-path construction in that neighborhood, while adding a functorial path object and, for the particular proof he chooses, bilateral units for path composition. The theorem's version~B instead uses a lifting property which is not part of Brown's basic definition. The comparison is between two ways of arranging enough structure, rather than an identification of their axiom lists.

The relevant later theorem is concrete. In \citet[Definition~1.1 and Proposition~1.8]{Cisinski2010}, derivable categories carry weak equivalences and fibrations with axioms arranged to include Brown's theory; the inclusion of the fibrant objects induces an equivalence on localizations. Theorem~2.21 constructs the associated weak derivator on finite direct categories, and Theorem~6.22 supplies diagram categories with model structures in Cisinski's sense under its stated completeness hypotheses. Those theorems explain the later organization of this material, but they do not imply that the two elementary axioms on p.~108 alone furnish a derivator on every small category.

One small warning concerns footnote~59. It says that the stability hypotheses on $\mathrm{Fib}$ will be used only in part~II and only in case~A. The proof of part~I in case~B has already used both base-change and composition stability on p.~93, to put $D\times_SD$ and its path object in $\underline{\mathrm{Fib}}\,S$. The stated theorem retains these hypotheses and is unaffected; the claim about where its proof uses them is inaccurate. This is a bookkeeping slip in the edition, not a missing assumption in the theorem.
\end{dcomment}

\begin{dcomment}{development}{p112}{Functorial replacement already gives the reflector}
Grothendieck constructs a path multiplication to prove the second triangle identity, but the comparison equivalence established in part~I already provides a shorter route in the functorial case. The simplification matters because it separates an explicit geometric model of the reflector from the hypotheses needed for its existence. The underlying equivalence is the standard deformation argument; compare \citet[Definition~2.2.1 and Remark~2.2.2, with arrows reversed]{Riehl2014}. We did not find the following two-localizer formulation stated in the sources checked for this package. We record it as a deduction from the supplied argument, without a claim of priority.

\textbf{Full replacement subcategories suffice.}\ours\quad
Let $\mathcal A$ be a category, $W$ a class containing the isomorphisms and satisfying two-out-of-three, and $\mathcal F\subseteq\mathcal A$ a full subcategory. Suppose that a functor $C:\mathcal A\to\mathcal F$ and a natural transformation $\eta:1_{\mathcal A}\to JC$ are given, where $J$ is the inclusion, with $\eta_X\in W$ for all $X$. Let $V\subseteq W$ be a second class such that
\[
V\cap\mathrm{Fl}(\mathcal F)=W\cap\mathrm{Fl}(\mathcal F).
\]
Assume the displayed localizations exist in the chosen universe. Then the canonical functor
\[
\rho:\mathcal A[V^{-1}]\longrightarrow\mathcal A[W^{-1}]
\]
has a fully faithful right adjoint $i$, represented on objects by $X\mapsto C(X)$. No path multiplication, lifting property, finite limits, or closure property of a class of fibrations is required.

\textit{Proof.}
Naturality says $C(f)\eta_X=\eta_Yf$ for every $f:X\to Y$, where inclusions of the full subcategory are suppressed. If $f\in W$, two-out-of-three gives $C(f)\in W$; since both endpoints lie in $\mathcal F$, this map also lies in $V$. Therefore $C$ descends to a functor
\[
\bar C:\mathcal A[W^{-1}]\longrightarrow\mathcal F[W^{-1}],
\]
and, by inclusion, defines $i:\mathcal A[W^{-1}]\to\mathcal A[V^{-1}]$. Write $\bar J:\mathcal F[W^{-1}]\to\mathcal A[W^{-1}]$ for the functor induced by $J$.

The components of $\eta$ give a natural isomorphism $1\to\bar J\bar C$ in $\mathcal A[W^{-1}]$. On restricting to $\mathcal F$, they also give a natural isomorphism $1\to\bar C\bar J$ in $\mathcal F[W^{-1}]$: fullness is what makes $\eta_F:F\to C(F)$ a morphism of that subcategory. Naturality extends from arrows to their inverses, so these are natural isomorphisms on the localizations themselves. Thus $\bar J$ is an equivalence, in particular faithful.

The same original transformation $\eta$ descends to $\widetilde\eta:1\to i\rho$ on $\mathcal A[V^{-1}]$, since $C$ sends $V\subseteq W$ into $V$ on $\mathcal F$. In the coarse localization define the counit $\varepsilon:\rho i\to1$ by $\varepsilon_X=[\eta_X]^{-1}$. It is natural there and invertible. The first triangle identity is
\[
\varepsilon_{\rho X}\,\rho(\widetilde\eta_X)=1_{\rho X}.
\]
For the second, apply naturality to $\eta_X$:
\[
C(\eta_X)\eta_X=\eta_{C(X)}\eta_X.
\]
After localization at $W$ we may cancel $\eta_X$. Both remaining arrows have source and target in $\mathcal F$. The faithfulness of $\bar J$ therefore gives
\[
[C(\eta_X)]=[\eta_{C(X)}]
\quad\text{already in }\mathcal F[W^{-1}].
\]
Its image in $\mathcal A[V^{-1}]$ yields
\[
i(\varepsilon_X)\widetilde\eta_{iX}
=[C(\eta_X)]^{-1}[\eta_{C(X)}]=1_{iX}.
\]
The triangle identities establish $\rho\dashv i$, and the invertible counit makes $i$ fully faithful. This proves the assertion.\hfill$\square$

Applied to Theorem-Scholie~2 with $\mathcal A=\mathcal M/S$, $\mathcal F=\underline{\mathrm{Fib}}\,S$, $V=W_S^0$, this proves part~II under version~A itself. In particular, the strengthening A$'$ is sufficient but unnecessary for the stated adjunction. There is no circular use of full faithfulness of $i$: the only faithfulness used in the proof is that of $\mathcal F[W^{-1}]\to\mathcal A[W^{-1}]$, already supplied by the two replacement isomorphisms in part~I. The fine-localization embedding is the conclusion.

The theorem also gives a precise version of the independence asserted in footnote~60. For any $F\in\mathcal F$ and any $X\in\mathcal A$, the canonical map
\[
\Hom_{\mathcal A[V^{-1}]}(X,F)\longrightarrow
\Hom_{\mathcal A[W^{-1}]}(X,F)
\]
is bijective: $\eta_F$ lies in $V$, so the adjunction formula with target $C(F)$ may be transported along this isomorphism. If $V\subseteq V'\subseteq W$ and both classes agree with $W$ on $\mathcal F$, then the intermediate localization $\mathcal A[V^{-1}]\to\mathcal A[(V')^{-1}]$ induces the unique bijection between these Hom sets which becomes the identity in the common coarse localization. This compatibility follows by composing the two displayed canonical bijections, so it involves no choice of a presentation by fractions.

A$'$ should nevertheless be read as a strengthening of A, exactly as the text says. Its short list on p.~112 does not restate Hypothesis~4. If one uses that list alone to \emph{construct} a weak-equivalence unit from a path object, the base-change argument on p.~102 must still be supplied. If A is retained, the unit's membership in $W$ is already assumed. The distinction avoids turning a compressed recap into a spurious counterexample to the theorem.
\end{dcomment}

\begin{dcomment}{exposition}{p115}{Which fibrations describe varying bases}
Grothendieck now returns from an abstract slice category to categories fibred over a category $S$. This changes the meaning of the word ``fibration.'' A Grothendieck fibration, used to encode a contravariant family of categories, is not by definition a fibration in a Quillen model structure on $\mathrm{Cat}$. Keeping these roles separate makes the passage to inverse images much easier to follow.

A strictly split Grothendieck fibration corresponds to a strict functor $F:S^{\mathrm{op}}\to\mathrm{Cat}$. Its total category has objects $(s,x)$ with $x\in F(s)$; an arrow $(s,x)\to(t,y)$ consists of $a:s\to t$ and a morphism $x\to F(a)y$ in $F(s)$. Given $g:T\to S$, the pullback of this total category has objects $(t,x)$ with $x\in F(g(t))$ and arrows described by the same rule. It is canonically the category associated to $F\circ g^{\mathrm{op}}$. Consequently an objectwise $W$-equivalence remains one after restriction, and the pullback functor passes to the relevant localizations.

For an arbitrary Grothendieck fibration one chooses cartesian lifts and obtains a pseudofunctor; a split replacement supplies a strict presentation. What the comparison theorem adds is that the localized answer does not depend on the splitting and can be computed using all functors over $S$ between fibred categories. The assertion concerns the localization of that category of functors. It does not say that a functor over $S$ preserves cartesian arrows, or that each resulting localized morphism is a strict natural transformation between the chosen strict diagrams. The replacement functor $\Phi$ is what reconciles these presentations.

For any functor $u:X\to Y$ over $S$, strict fibres of a pullback satisfy
\[
(T\times_SX)_t\cong X_{g(t)},\qquad
(T\times_SY)_t\cong Y_{g(t)}.
\]
Thus fibrewise weak equivalences are preserved by arbitrary base change. This proves the displayed inclusion on p.~117 without invoking properness. The comparison with $\HOT_W(S)$ is what allows this elementary operation on fibred categories to compute the inverse image of a relative homotopy type. The corrected target localizer is $W_{T}^{\mathrm f}$, not the printed $W_S^{\mathrm f}$. The next pages explain why applying the same ordinary pullback to every object of $\mathrm{Cat}/S$, without a suitable replacement, need not have this property for the directional localizer $W_S^{\mathrm d}$.

The reference to Thomason on p.~114 concerns a different existence theorem. \citet[Theorem~4.9 and Corollary~5.5]{Thomason1980} constructs a proper model structure on small categories whose weak equivalences are detected by the nerve and whose fibrations are detected by $\mathrm{Ex}^2N$. That supplies model-category factorizations for the ordinary homotopy theory of categories. It does not identify those fibrations with Grothendieck fibrations, and it does not by itself verify the same factorization statement for every other localizer $W$. The geometric presentations of diagrams remain part of the argument even when a model structure is available.
\end{dcomment}

% Commentary on manuscript pp.118--145.

\begin{dcomment}{exposition}{p118}{Computing restriction on a suitable representative}
Grothendieck wants restriction of relative homotopy types to retain the simplicity of a fibre product. The work on the several categories of models now has a precise use: one changes the representative before taking the fibre product, so that a calculation in ordinary categories gives the derived answer. Throughout these comments we retain his presheaf convention. Thus a diagram over $S$ varies contravariantly in $S$, and the value represented by a category $p:X\to S$ at an object $s$ is the homotopy type of the comma category
\[
 s\backslash X=(s\downarrow p).
\]
For $W=W_\infty$, this means its nerve, regarded as a space. In particular, $W_S^{\mathrm d}$ consists of the maps which induce weak equivalences on all these comma categories. The strict fibre $X_s$ is a different category.

The fibration replacement $\Phi_SX$ has fibre $s\backslash X$. The natural map $X\to\Phi_SX$ belongs to $W_S^{\mathrm d}$. If $X\to S$ is smooth in the sense being used here, the comparison $X_s\to s\backslash X$ is a weak equivalence, and smoothness is preserved by change of base. Consequently the strict pullback of a weak equivalence between smooth representatives is again a weak equivalence of the relevant relative kind. This proves the asserted formula
\[
 f^*[X]\simeq[X\times_S S']
\]
on smooth representatives, and proves its compatibility with maps between them. A comparison equivalence from the smooth models to the whole localized category then supplies the formula for arbitrary objects after applying $\Phi_S$. This last use of the comparison theorem is essential: a formula on objects alone would not define a functor on the localized category.

Here is the elementary obstruction which makes the replacement necessary. Put $S=[1]=(0\to1)$ and let $u:\{1\}\to S$ be the terminal-object inclusion, viewed as a map over $S$. For each $s=0,1$, both $s\backslash\{1\}$ and $s\backslash S$ are contractible; hence $u\in W_S^{\mathrm d}$. Pulling back along $\{0\}\to S$, however, gives $\varnothing\to e$, which is not in $W_\infty$. Thus the fibre product does not descend on all of $\Cat/S$. The parenthetical suggestion that every base other than the point causes this problem should not be read as a theorem: if $S$ is discrete, $s\backslash X=X_s$, and every change of base preserves this fibrewise condition. What fails is the unrestricted assertion for general bases.

The preceding comparison of smoothness with genuine fibrations has its systematic treatment in \citet{Maltsiniotis2005Asphericite}; its purpose is precisely to retain the homotopical consequence of cartesian lifting while allowing more models. This extends the earlier companion's discussion of local isomorphisms and cofinality in DER\_I-a, ms.~pp.~52--66, rather than changing the direction of the comma category.
\end{dcomment}

\begin{dcomment}{examples}{p120}{The two replacements do not compute the same inverse image}
Grothendieck now offers two ways to improve $X\to S$, but they have different effects on the relative homotopy type. The second isomorphism in Corollary~2 is a genuine error in the edition, not a fault of the supplied transcription. The distinction was already explicit on manuscript p.~66: $X\to C(X)$ belongs to $W_S$, and need not belong to $W_S^{\mathrm d}$. The object $C(X)$ represents the locally constant reflection of $[X]$, whereas $\Phi_SX$ represents $[X]$ itself.

To see the error without any machinery, take $W=W_\infty$, $S=[1]$, and $X=\{0\}$ over $S$. This is a split fibration: its presheaf of discrete fibres is
\[
 F(1)=\varnothing\longrightarrow F(0)=e.
\]
For the inclusion $j:\{1\}\to S$, the correct restriction is therefore $j^*[X]=\varnothing$. On the other hand, the map of total categories $X\to S$ is a weak equivalence, since both nerves are contractible. Hence $\rho_S[X]=\rho_S[S]$. The relative type represented by $S\to S$ is the constant point, and the fully faithful embedding $i_S$ identifies $[C(X)]$ with this constant point. It follows that
\[
 j^*[C(X)]\simeq e\not\simeq\varnothing=j^*[X].
\]
All the categories in this example are finite; no issue of size, saturation, or choices of paths is involved.

The two correct formulas, with the ambient categories indicated, are
\[
 \begin{aligned}
 f^*[X]&\simeq[\Phi_SX\times_S S']&&\text{in }\HOT_W(S'),\\
 i_{S'}f^*_{\mathrm{lc}}\rho_S[X]
   &\simeq[C(X)\times_S S']&&\text{in }\HOT_W(S').
 \end{aligned}
\]
The second uses the compatibility $f^*i_S\simeq i_{S'}f^*_{\mathrm{lc}}$, which holds because restriction preserves locally constant objects. It does not assert that restriction commutes with the reflector $\rho$. The latter assertion is the separate problem of Proposition~2. If $[X]$ is already locally constant, the unit $[X]\to i_S\rho_S[X]$ is invertible, and the formulas agree. For arbitrary $X$ one may also apply $\rho_{S'}$ to the first formula, but identifying the result with the second requires the additional commutation property being studied below.

This example also explains the later correction on manuscript p.~127. There Grothendieck does use $C(X)$ to compute the restriction of $\rho_S[X]$, which is the appropriate operation. The diagram of fifteen categories does not license interchanging its different localizers; it gives several presentations of each of three distinct localized categories.
\end{dcomment}

\begin{dcomment}{exposition}{p122:déjà connu ici. Cf. théorème}{Proper base change and the homology of comma categories}
Grothendieck's appeal to relative homology is intended to turn a local test on categories into an isomorphism in a single category of coefficients. Let $u:X\to Y$ lie over $S$, let $p_X,p_Y$ be the structure maps, and let $a\in\mathbf D(e)$. Pointwise left Kan extension identifies the value at $s$ of the natural map
\[
 (p_X)_!a_X\longrightarrow(p_Y)_!a_Y
\]
with the map on homology of $s\backslash X\to s\backslash Y$ with constant coefficient $a$. Since evaluations are conservative, these relative homology maps are invertible for every $a$ exactly when all the comma-category maps are $\mathbf D$-equivalences. This is the argument hidden in the reference to Der~4~bis. The word ``homology'' here denotes an object of the coefficient category, not merely its homotopy groups or ordinary integral homology.

Now let $g:T\to S$ be $\mathbf D$-proper. Write $X_T=X\times_ST$ and $Y_T=Y\times_ST$. The defining base-change isomorphisms give a natural commutative square
\[
\begin{tikzcd}[column sep=large]
g^*(p_X)_!a_X\arrow[r]\arrow[d,"\sim"']&g^*(p_Y)_!a_Y\arrow[d,"\sim"]\\
(p_{X_T})_!a_{X_T}\arrow[r]&(p_{Y_T})_!a_{Y_T}.
\end{tikzcd}
\]
If the upper map is invertible, so is the lower one. Applying the preceding pointwise criterion over $T$ proves the lemma. This supplies the whole proof and also explains its naturality in $a$ and in $u$. The four families in the following list give equivalent tests because they contain the elementary comma projections and because proper base change preserves the condition. The phrase ``three types'' in the text is an uncorrected counting slip.

There is no unresolved conjecture here for the ordinary homotopy derivator. Grothendieck supplies his own categorical proof later in this chapter, manuscript pp.~168--169. For $t\in T$ and $s=g(t)$, properness makes
\[
 t\backslash T\longrightarrow s\backslash S
\]
a universal $W$-equivalence. Pullback along a universal $W$-equivalence preserves $W$: in the pullback square for $v:A\to B$, both projections to $A$ and $B$ lie in $W$, so two-out-of-three identifies the conditions $v\in W$ and $v'\in W$. Apply this to $s\backslash u$. This later argument proves the conjectured preservation of $W_S^{\mathrm d}$ directly, without first constructing every derivator operation. We have checked the continuation in the full edition, separately from the attached excerpt \citep[VI, ms.~pp.~168--169]{Grothendieck1990VI}.

The modern organization is the theory of homotopy exact squares. \citet[Proposition~2.13 and Theorem~4.32]{Maltsiniotis2012Exact} separates the categorical criterion for properness from the criterion for a derivator's base-change map to be invertible. These results are a systematic later formulation of the mechanism used here. The tentative converse on p.~123 also holds for $W_\infty$, as the following elementary test proves. Its proof matters: preserving a class of maps and universal exactness of squares should not simply be identified without establishing their comparison.
\end{dcomment}

\begin{dcomment}{development}{p123:Ce qui devrait être le cas si (et seulement si?)}{Testing the converse on terminal objects of slices}
\textbf{Proposition (the ordinary properness criterion).}\ours\
For a functor $f:T\to S$, strict pullback preserves $W_S^{\mathrm d}$-equivalences, with $W=W_\infty$, if and only if $f$ is $W_\infty$-proper.

\textit{Proof.} Properness implies preservation by the argument discussed above. For the converse, fix $s\in S$, let $e_s\to S$ select $s$, and consider
\[
 u_s:e_s\longrightarrow S/s,
\]
where $u_s$ selects the terminal object $(s,\mathrm{id}_s)$ and $S/s\to S$ forgets the arrow to $s$. For each $r\in S$, the source $r\backslash e_s$ is the discrete category $\mathrm{Hom}_S(r,s)$. The target $r\backslash(S/s)$ has objects factorizations $r\to v\to s$ and decomposes as a disjoint union indexed by their composites $\gamma:r\to s$. The component indexed by $\gamma$ has terminal object $r\xrightarrow{\gamma}s\xrightarrow{\mathrm{id}}s$. Thus $r\backslash u_s$ is a weak equivalence for every $r$, proving $u_s\in W_S^{\mathrm d}$.

Its pullback along $f$ is the fibre inclusion
\[
 j_s:T_s\longrightarrow T/s
\]
over $T$, where $T/s=(f\downarrow s)$. By hypothesis, $t\backslash j_s$ is a weak equivalence for every $t\in T$. The target $t\backslash(T/s)$ decomposes into components indexed by arrows $\alpha:f(t)\to s$: an object is a pair $h:t\to v$, $\beta:f(v)\to s$ with composite $\beta f(h)=\alpha$. Each component has initial object $(\mathrm{id}_t,\alpha)$. The source $t\backslash T_s$ decomposes using the same labels, assigning $h:t\to v$, $f(v)=s$, to $f(h)=\alpha$. Its $\alpha$-component is exactly
\[
 (t,\alpha)\backslash j_s.
\]
Since a weak homotopy equivalence of these disjoint unions induces a bijection on connected components and equivalences on their homotopy groups, each such source component is nonempty and weakly contractible. Thus $j_s$ is $W_\infty$-coaspherical for every $s$. The usual categorical properness criterion says precisely that these fibre inclusions are coaspherical; see the proof of \citet[Proposition~2.13]{Maltsiniotis2012Exact}. Hence $f$ is $W_\infty$-proper.\hfill$\square$

This settles the parenthetical ``if and only if?'' in the ordinary case. The proof deliberately states that scope: extraction of the individual summands uses detection of connected components by $W_\infty$. A statement for an arbitrary localizer must retain the hypotheses which allow that step. The calculation is supplied here as an explicit application of the established properness criterion, without a priority claim.
\end{dcomment}

\begin{dcomment}{status}{p126}{Saturation completes the equivalence in Proposition 2}
Grothendieck catches the precise point at which his attempted converse uses more than the universal property of localization. A functor descending to a localization sends the chosen weak equivalences to isomorphisms; without saturation this says nothing about whether their representatives were members of the originally chosen class. His footnote is therefore a valid qualification, not an unnecessary doubt.

In the ordinary case the missing step is short. Suppose that $W$ is strongly saturated in $\Cat$: a map belongs to $W$ whenever its image in $\Cat[W^{-1}]$ is invertible. Let $\mathcal F$ be any subcategory of $\Cat/S$, and put $W_{\mathcal F}=W_S\cap\mathrm{Mor}(\mathcal F)$. The forgetful functor induces
\[
 \mathcal F[W_{\mathcal F}^{-1}]\longrightarrow\Cat[W^{-1}].
\]
If $u\in\mathrm{Mor}(\mathcal F)$ becomes invertible on the left, its underlying functor becomes invertible on the right. Strong saturation of $W$ implies $u\in W_{\mathcal F}$. Thus strong saturation passes to every such restricted class; fullness of $\mathcal F$ is not needed for this observation.

In the ordinary case $W=W_\infty$, this historical qualification leaves no surviving gap: the hypothesis holds. The nerve identifies its localization with the homotopy theory of spaces, and a map of simplicial sets invertible in the homotopy category is a weak equivalence. Applying the observation to each target category $\mathcal F'$ supplies the missing implication in Proposition~2. More explicitly, if
\[
 \rho_{S'}f^*\simeq f^*_{\mathrm{lc}}\rho_S,
\]
then a coarse weak equivalence $u$ between smooth, fibred, or split-fibred models is sent to an isomorphism after target localization; the observation puts its strict pullback back in $W_{S'}$. Conversely, any one of these preservation conditions produces the two commutative localization squares written in the text. Restriction on the locally constant subcategories identifies the resulting right-hand functor with $f^*_{\mathrm{lc}}$. This proves all four conditions equivalent under ambient strong saturation.

The obstruction is visible in the example just given. For $j:\{1\}\to[1]$ and the diagram $\varnothing\to e$, one has
\[
 \rho_ej^*F=\varnothing,
 \qquad j^*_{\mathrm{lc}}\rho_{[1]}F=e.
\]
Thus restriction preserves the locally constant subcategory while failing to commute with its reflection. The two assertions must remain distinct. Under the chapter's axiom L6~bis, a Hot-fibration gives a sufficient condition for commutation, as the corollary states. For a general $W$ we retain the stated dependence on its axioms; the preceding proof does not silently replace $W$ by its saturation. Nor do the unqualified closure axioms alone identify $W$ with $W_\infty$: the maximal class of all functors is an immediate counterexample to such an inference. An identification theorem must retain its nondegeneracy restriction, such as $W\subseteq W_0$, excluding that maximal localizer.
\end{dcomment}

\begin{dcomment}{development}{p126}{A Hot-fibration which does not preserve colocal equivalences}
Grothendieck explicitly asks for a counterexample separating the two pullback properties. A small category with a deformation to its base supplies one. The construction is useful because every coarse homotopy calculation is trivial, while a single comma category retains a nontrivial set of lifts. We checked the treatments of smoothness and exact squares in \citet{Maltsiniotis2005Asphericite,Maltsiniotis2012Exact}; we did not find this particular finite presentation there. The result below is a calculation following the manuscript's question, without a claim of priority.

\textbf{Proposition (a finite separation example).}\ours\
There is a universal $W_\infty$-equivalence $p:A\to[1]$, hence a Hot-fibration, for which the strict pullback functor does not preserve $W^{\mathrm d}$.

\textit{Proof.} Define $A$ by objects $r,x,z,t$, generating arrows
\[
 r\xrightarrow{a}x\xrightarrow{b}z,
 \qquad x\overset{\alpha}{\underset{\beta}{\rightrightarrows}}t,
 \qquad \alpha a=\beta a,
\]
and no further generating arrows or relations. Thus $r$ is initial, the full subcategory on $r,x,z$ is the chain $r\to x\to z$, and there is no arrow $z\to t$. Let $p$ take $r,x,z$ to $0$ and $t$ to $1$. The two arrows $\alpha,\beta$ have image $0\to1$. Define a section $k:[1]\to A$ by $k(0)=r$, $k(1)=t$, and $k(0\to1)=\alpha a=\beta a$. There is a natural transformation
\[
 kp\Longrightarrow\mathrm{id}_A
\]
whose components at $r,x,z,t$ are respectively $\mathrm{id}_r,a,ba,\mathrm{id}_t$. Naturality at $\alpha$ and $\beta$ is exactly the imposed relation. All its components map to identities in $[1]$, so this is a transformation over the base.

After any base change $T\to[1]$, both the section and this transformation pull back. The projection $A\times_{[1]}T\to T$ consequently has a section whose composite is naturally homotopic to the identity. Its nerve is a homotopy equivalence. This proves that $p$ is a universal $W_\infty$-equivalence.

To verify explicitly the Hot-fibration assertion in the terminology of these pages, consider any such pullback $A_T\to T$ and $s\in T$. Its strict fibre, $s\backslash A_T$, and $A_T/s$ all have contractible nerves. For the comma categories this follows by pulling back the universal equivalence along $s\backslash T\to T$ and $T/s\to T$ respectively; the bases have initial or terminal objects. The strict fibre is obtained by pullback to $e$. The canonical inclusions from the fibre to the two comma categories are therefore weak equivalences, as are all the transition functors between these contractible comma categories. This remains true after every further base change, which gives the required Hot-fibration condition in either of the formulations recalled in the chapter.

Now let $u:e=\{1\}\to[1]$. Both comma categories $0\backslash e$ and $1\backslash e$ are points, whereas the corresponding undercategories of $[1]$ are contractible. Hence $u\in W_{[1]}^{\mathrm d}$. Its strict pullback is the inclusion $v:\{t\}\to A$. The category $x\backslash\{t\}$ is discrete with two objects, namely $\alpha$ and $\beta$: a morphism between them would require their equality. In contrast, $x\backslash A$ has initial object $\mathrm{id}_x$. Thus
\[
 x\backslash\{t\}\longrightarrow x\backslash A
\]
is not a weak equivalence, since it does not induce a bijection on connected components. Therefore $v\notin W_A^{\mathrm d}$.\hfill$\square$

The same construction confirms the distinction at the level of coarse equivalences. Since $p$ is universal, for every map $X\to Y$ over $[1]$ the projections from its pullback to $X$ and $Y$ are weak equivalences. Two-out-of-three proves preservation of coarse $W$. Failure of preservation of $W^{\mathrm d}$ is thus compatible with the corollary's commutation of restriction with locally constant reflection: that corollary computes ordinary restriction on suitable smooth models, not by raw pullback of every representative.
\end{dcomment}

\begin{dcomment}{development}{p128}{Locally constant coefficients have adjoints but fail the pointwise axiom}
Grothendieck asks whether the locally constant categories themselves form a derivator, and suggests looking for an obstruction to the existence of adjoints. In the ordinary theory the obstruction occurs elsewhere. Restriction between categories of locally constant coefficients does have both adjoints: this is \citet[Corollary~3.10]{Cisinski2009LocallyConstant}. What fails, when the base category itself is the diagram variable, is the pointwise formula for Kan extensions. We spell out the distinction because otherwise a later theorem about locally constant derivators appears to contradict the answer.

\textbf{Proposition (the walking-arrow obstruction).}\ours\
Let $\mathcal L(S)$ be the full subcategory of $\operatorname{Ho}(\mathrm{sSet}^{S^{\mathrm{op}}})$ consisting of diagrams which send every arrow of $S$ to a weak equivalence. With its usual restriction functors, $S\mapsto\mathcal L(S)$ does not satisfy the left pointwise Kan-extension axiom on all small categories.

\textit{Proof.} Let $S=[1]$ and let $j:e\to S$ select $0$. Evaluation $j^*:\mathcal L(S)\to\operatorname{Ho}(\mathrm{sSet})$ is an equivalence. Indeed, a representative is a weak equivalence $F(1)\to F(0)$, and the natural map from this diagram to the constant diagram with value $F(0)$ is an objectwise weak equivalence. The constant-diagram functor is consequently a quasi-inverse to evaluation; the same natural maps prove the statement on morphisms, not just on isomorphism classes of objects. Any left adjoint $j_!^{\mathcal L}$ to this equivalence is isomorphic to that quasi-inverse. Hence
\[
 (j_!^{\mathcal L}e)(1)\simeq e.
\]
In the presheaf convention, the pointwise left Kan-extension formula at $1$ uses the comma category $(1\downarrow j)$. This category is empty because there is no arrow $1\to0$. The formula would therefore give
\[
 (j_!^{\mathcal L}e)(1)\simeq
 \operatorname{hocolim}_{(1\downarrow j)} e=\varnothing,
\]
which is impossible. The empty indexing category has its unique diagram, and its left extension to a point is necessarily the initial space, so no alternate choice of adjoint changes this conclusion.\hfill$\square$

This proof develops the question asked here; we make no claim of priority for the obstruction. We checked the discussion of locally constant coefficients in \citet[Sections~1--3]{Cisinski2009LocallyConstant}; we did not find this particular walking-arrow formulation there. The distinction with his Theorem~3.12 is precise. Fixing a category $A$, he constructs the derivator of diagrams on $A\times I^{\mathrm{op}}$ which are locally constant in the $A$ direction, while $I$ remains the unrestricted diagram variable. That is a derivator. The construction considered by Grothendieck here instead makes the entire varying base locally constant. It removes exactly the noninvertible diagram arrows needed to make pointwise extension by the empty homotopy type possible.

For ordinary spaces there is also a useful interpretation of why the adjoints nevertheless exist. A locally constant diagram factors homotopically through the groupoid completion of its base, and its restriction is restriction of a family over the induced map of spaces. Left and right extension over that map exist. Their fibres are calculated by homotopy fibres of the map of groupoid completions, which need not agree with the nerves of the ordinary comma categories occurring in Der~4. In the example the groupoid completion of $[1]$ is contractible, so extending the point remains the point everywhere. This is exactly the calculation above.
\end{dcomment}

\begin{dcomment}{exposition}{p130:Donc on ne peut espérer trouver une description}{Natural transformations and cartesian transport}
Grothendieck wants the two-dimensional structure to be visible on the same geometric models used for restriction. A natural transformation $u:f\Rightarrow g$, with $f,g:T\to S$, has to induce $u^*:g^*\Rightarrow f^*$ because the diagrams are presheaves. This direction is forced by the contravariance of the coefficient system; it is not a separate duality convention to be chosen later.

For a split fibration classified by $F:S^{\mathrm{op}}\to\Cat$, the component over $t\in T$ is simply
\[
 F(u_t):F(g(t))\longrightarrow F(f(t)).
\]
Naturality of $u$ makes these functors a map of $T$-diagrams. Identity transformations and vertical and horizontal composition give the required identities by functoriality of $F$. Localizing at objectwise weak equivalences therefore retains the whole two-functor, not just its action on objects and one-morphisms.

For an ordinary Grothendieck fibration, choose a cleavage. Lift $u_t:f(t)\to g(t)$ cartesianly to obtain a reindexing functor from the fibre over $g(t)$ to the fibre over $f(t)$. The cartesian universal property gives the comparison maps for arrows of $T$. Two choices of cartesian lifts have a unique vertical isomorphism compatible with their maps to the same target. This uniqueness makes the comparison coherent under composition; different cleavages therefore give canonically isomorphic maps after localization. A strict choice-free map of total categories is more than the argument requires. Split fibrations provide a convenient strict presentation, and ordinary fibrations recover it up to these canonical comparisons.

Smoothness or Hot-fibrancy controls homotopy types of fibres but does not provide a specified cartesian lifting operation. It therefore does not, by itself, produce the strict functor demanded in the displayed question. One may still calculate the induced transformation in the localized category by replacing the model with $\Phi_SX$. This is the useful content of the negative observation on p.~129; it should not be turned into the assertion that every individual nonfibred example admits no fibre-to-fibre functor at all.

The passage concerning Der~2 must be read with its cancellation. Strong saturation of $W$ tests whether an actual map of the coefficient category becomes invertible. Der~2 tests an arbitrary morphism already present in a localized diagram category, which may be represented by a zigzag, and requires objectwise detection of its invertibility. These are different levels of assertion. For the standard space derivator the latter property holds, as does the full derivator structure; see \citet[Proposition~1.30]{Groth2013} for the modern axioms and model-category examples. The text does not prove it for every bare pair $(\mathcal M,W)$, and its cancelled assertion should not be restored on the strength of saturation alone.
\end{dcomment}

\begin{dcomment}{development}{p130:Pour les axiomes essentiels sur les dérivateurs}{Strong saturation does not imply evaluation conservativity}
The cancellation in footnote~66 concerns a real distinction. Strong saturation of $W$ does imply strong saturation of the pointwise class $W^I$ in $\mathcal C^I$: if an actual natural transformation becomes invertible after diagram localization, evaluation makes each of its components invertible in $\mathcal C[W^{-1}]$, hence a member of $W$. But Der~2 concerns every arrow of the localized diagram category. Such an arrow may be a zigzag, and saturation alone does not replace it by an actual natural transformation. The following finite example gives a negative answer to the unrestricted question. We did not locate this construction in the sources checked; we provide its complete proof without a claim of priority.

\textbf{A strongly saturated relative category failing Der~2.}\ours\quad
There is a category $\mathcal C$ with four objects and fourteen morphisms, and a strongly saturated class $W$, such that
\[
\mathcal C^{[1]}[(W^{[1]})^{-1}]
\longrightarrow
\bigl(\mathcal C[W^{-1}]\bigr)^{[1]}
\]
does not reflect isomorphisms. In fact a noninvertible endomorphism has identity evaluations at both vertices.

\textit{Proof.}
Let $\mathcal C$ have objects $K,A,B,Q$ and generators
\[
a:K\to A,\qquad h:A\to B,\qquad g:B\to A,
\qquad k:K\to B,\qquad w:B\to Q,
\]
subject only to
\[
gh=1_A,\qquad gk=a,\qquad wha=wk.
\]
Put $u=ha$, $v=k$, $t=wk=wha$, and $e=hg$. The Hom sets, with source in the row and target in the column, are
\[
\begin{array}{c|cccc}
 &K&A&B&Q\\ \hline
K&\{1_K\}&\{a\}&\{u,v\}&\{t\}\\
A&\varnothing&\{1_A\}&\{h\}&\{wh\}\\
B&\varnothing&\{g\}&\{1_B,e\}&\{w,we\}\\
Q&\varnothing&\varnothing&\varnothing&\{1_Q\}.
\end{array}
\]
Indeed, the three relations reduce every path to a displayed representative. The three potentially distinct parallel pairs really are distinct. To see this, map the presentation to sets by taking $K=A=\{*\}$, $B=\{0,x,y\}$, $Q=\{p,q\}$, with $h(*)=0$, $g$ constant, $a$ the unique map, $k(*)=x$, and $w(0)=w(x)=p$, $w(y)=q$. This model distinguishes $u$ from $v$, $1_B$ from $e$, and $w$ from $we$. Thus the table is exact.

Let $W$ consist of the four identities and $w$. To verify strong saturation, use a different functor $F:\mathcal C\to\mathrm{Set}$:
\[
F(K)=\{*\},\quad F(A)=\{0,1\},\quad
F(B)=F(Q)=\{0,1,2\}.
\]
Take $F(a)(*)=F(k)(*)=0$, let $F(h)$ be the inclusion, let $F(g)$ fix $0,1$ and send $2$ to $0$, and put $F(w)=1$. These formulas respect the relations. The Hom table shows that the inverse image under $F$ of the bijections is exactly $W$. Since $F$ factors through $\mathcal C[W^{-1}]$, an actual arrow inverted by that localization must belong to $W$. The converse is built into localization, proving strong saturation.

Consider now the full subcategory $\mathcal R\subseteq\mathcal C^{[1]}$ on the four arrows
\[
a:K\to A,\qquad u:K\to B,\qquad
v:K\to B,\qquad t:K\to Q.
\]
Every square between them has source component $1_K$, so denote such a square by its target component, adding its source-object subscript when necessary. A direct use of the Hom table shows that $\mathcal R$ is generated by
\[
h:a\to u,\quad g_u:u\to a,\quad g_v:v\to a,
\quad w_u:u\to t,\quad w_v:v\to t,
\]
with the sole relation $g_uh=1_a$. For precision, its other nonidentity arrows are
\[
e_u=hg_u,\quad e_v=hg_v,\quad
wh=w_uh,\quad (we)_u=w_uhg_u,\quad (we)_v=w_uhg_v.
\]
There is no arrow $a\to v$ or $u\to v$. These formulas, together with the identities, list all fourteen arrows of $\mathcal R$. The only nonidentity pointwise $W$-maps in $\mathcal R$ are $w_u$ and $w_v$.

Define a functor $H:\mathcal R\to\mathrm{Set}$ by assigning the same two-element set $E=\{0,1\}$ to each object, assigning the identity of $E$ to $h,g_u,w_u,w_v$, and assigning the constant map with value $0$ to $g_v$. The generating relation is respected, so this defines a functor. It sends both pointwise weak equivalences to identities.

One must ensure that a witness on this full subcategory remains valid after localizing the whole arrow category. Here it extends directly. Give $K$ degree $0$ and $A,B,Q$ degree $1$. There is no arrow of $\mathcal C$ which decreases degree. The source and target degrees of an object of $\mathcal C^{[1]}$ can therefore be only
\[
(0,0)<(0,1)<(1,1).
\]
The middle fiber is exactly $\mathcal R$, and the lower fiber consists only of $1_K$. Extend $H$ to a functor $\widetilde H:\mathcal C^{[1]}\to\mathrm{Set}$ by assigning $\varnothing$ to $1_K$ and a singleton to every object in the upper fiber. Within the upper fiber and between different fibers use the unique possible maps. Since natural transformations cannot decrease either degree, these prescriptions define a functor. Every pointwise $W$-map preserves degree; in the middle fiber it is an identity or one of the two maps already considered, and in the other fibers its image is an identity. Thus $\widetilde H$ inverts all of $W^{[1]}$ and descends to the full localized arrow category.

Finally form
\[
\alpha=[g_v][w_v]^{-1}[w_u][h]:a\longrightarrow a.
\]
Its source evaluation is $1_K$, and its target evaluation is
\[
[g][w]^{-1}[w][h]=[gh]=1_A.
\]
But its image under the descended $\widetilde H$ is the constant endomorphism of $E$, which is not bijective. Therefore $\alpha$ is not invertible. This proves the assertion. In the contravariant convention of the manuscript one takes $S=[1]^{\mathrm{op}}$.\hfill$\square$

The example identifies exactly what a fraction theorem supplies in the positive cases. Suppose every localized diagram arrow can be written
\[
\alpha=[v]^{-1}[f][u]^{-1},
\qquad X\xleftarrow{u}X'\xrightarrow{f}Y'
\xleftarrow{v}Y,
\]
with $u,v$ pointwise in $W$. If each evaluation of $\alpha$ is invertible, so is every $[f_i]=[v_i]\alpha_i[u_i]$. Strong saturation makes $f$ pointwise weak, and hence $\alpha$ invertible. A three-arrow calculus gives this representation; compare \citet[Proposition~4.1.11(ii)]{Low2013HomotopicalAlgebra}. Grothendieck later obtains Der~2 with cofibration structures on the diagram categories and strong saturation, in Chapter~XIX, Corollary~2.4 and the explanation on manuscript p.~23 \citep{Grothendieck1990XIX}. For the published comparison with coherent diagrams, \citet[Theorem~3.33]{Lenz2018} identifies the prederivator of a cofibration category with that of its quasicategory of frames. These positive results require structure absent from the finite counterexample; they settle the intended homotopical cases without establishing the unrestricted assertion cancelled here.
\end{dcomment}

\begin{dcomment}{development}{p132}{The uniform bound in an infinite product is a real obstruction}
Grothendieck isolates the correct combinatorial issue: every morphism in an ordinary localization is represented by a finite zigzag, and a zigzag in a product has a single finite length shared by all coordinates. Coordinatewise existence of zigzags does not supply a common bound. The general question on p.~131 has a negative answer even when every chosen class of weak equivalences is strongly saturated. The special homotopy theory of spaces has additional structure which avoids the obstruction.

\textbf{An infinite-product counterexample.}\ours\quad For each integer $n\geq1$, let $P_n$ be the finite poset whose Hasse diagram is the alternating interval
\[
 0\longrightarrow1\longleftarrow2\longrightarrow3
   \longleftarrow\cdots\longrightarrow2n-1\longleftarrow2n.
\]
There are no composable pairs of nonidentity arrows. Set $W_n=\mathrm{Mor}(P_n)$. Its localization is the groupoid completion of a tree. There is exactly one morphism between any two objects: after deleting adjacent inverse pairs, the unique undirected interval between the two vertices is the only remaining path. Thus each $P_n[W_n^{-1}]$ is equivalent to $e$.

Put $P=\prod_{n\geq1}P_n$ and $W=\prod_nW_n=\mathrm{Mor}(P)$. Consider the objects $x=(0)_n$ and $y=(2n)_n$. Every arrow of $P$ changes its $n$th coordinate by at most one. Traversing an arrow backwards after localization has the same bound. Consequently a zigzag of length $r$ from $x$ can end only at an object $z$ with $|z_n|\leq r$ for every $n$. Since the distances $2n$ are unbounded, there is no finite zigzag from $x$ to $y$. Hence
\[
 \mathrm{Hom}_{P[W^{-1}]}(x,y)=\varnothing,
 \qquad
 \mathrm{Hom}_{\prod_nP_n[W_n^{-1}]}(x,y)=\{*\}.
\]
The canonical comparison is not full. All the classes $W_n$, and also $W$, are strongly saturated because they contain every morphism. This example shows that saturation cannot repair the infinite-product problem.

For a finite product, there is no such obstruction. Here is a proof that also handles relations between zigzags. Let $(\mathcal C,W)$ and $(\mathcal E,V)$ be relative categories with identities in their designated classes. A functor $H:\mathcal C\times\mathcal E\to\mathcal B$ inverts $W\times V$ if and only if it inverts each $(w,\mathrm{id})$ and $(\mathrm{id},v)$. For every $c$, the functor $H(c,-)$ therefore factors through $\mathcal E[V^{-1}]$. Natural transformations between these functors factor uniquely as well, by the fully faithful part of the universal property of localization. We obtain a functor
\[
 \mathcal C\longrightarrow
 \operatorname{Fun}(\mathcal E[V^{-1}],\mathcal B)
\]
which takes $W$ to natural isomorphisms. Localizing once more gives a functor from $\mathcal C[W^{-1}]\times\mathcal E[V^{-1}]$ to $\mathcal B$, unique with the same natural-transformation property. Thus this product has the universal property of $(\mathcal C\times\mathcal E)[(W\times V)^{-1}]$. Induction handles any finite family. This argument avoids assuming an unjustified strict identity between two chosen categories of paths.

The desired infinite coproduct axiom nevertheless holds for the ordinary homotopy derivator. For example, simplicial-set diagram categories over $\coprod_iS_i$ identify with products of diagram categories over the $S_i$, with their objectwise weak equivalences and their projective model structures. The product model structure is coordinatewise. Choose cofibrant and fibrant replacements coordinatewise; maps between cofibrant-fibrant objects represent morphisms in the homotopy category, and two such maps agree there exactly when they are homotopic. Cylinder objects and the witnesses of homotopy can also be chosen coordinatewise. Therefore
\[
 \operatorname{Ho}\Bigl(\prod_i\mathcal M_i\Bigr)
   \simeq\prod_i\operatorname{Ho}(\mathcal M_i)
\]
for a set-indexed family of model categories. The usual universe convention and the ordinary axiom of choice are understood. Applying this to the diagram categories proves Der~1, including arbitrary small coproducts, for $\HOT_{W_\infty}$. The same argument does not establish it for an arbitrary pair $(\mathcal M,W)$ lacking replacement and homotopy calculus.

The counterexample and the proof above are supplied here as worked calculations of the question in the manuscript, without a priority claim. The modern positive conclusion for model-derived derivators is standard \citep[Proposition~1.30]{Groth2013}.
\end{dcomment}

\begin{dcomment}{context}{p134}{Integration of a constant family}
Grothendieck is moving from a presentation of diagram categories to their operations. The word ``integration'' expresses the geometric idea: a family of spaces parametrized by $S$ has a total space, and the left adjoint to restriction should recover that total object. The source's names Der~3 and Der~4 are its historical numbering; in the commonly used modern list, existence of Kan extensions is Der~3 and their pointwise calculation is Der~4 \citep[Definition~1.5]{Groth2013}. Renumbering the manuscript silently would obscure the argument.

For a constant family the expected calculation can be carried out at once. Take an actual category $A$ representing the coefficient $a\in\mathrm{Hot}$ and the projection $S\times A\to S$. It is a split fibration with constant fibre $A$. If integration is induced by composition of structure maps, the pushforward along $p:S\to e$ is the category $S\times A$. The nerve preserves products degreewise:
\[
 N(S\times A)=NS\times NA.
\]
Consequently
\[
 p_!p^*a\simeq B(S)\times a
\]
in the homotopy category of spaces, where $B(S)$ denotes the homotopy type of $NS$. This is a product of homotopy types. For coefficients in a general cocomplete homotopy theory, the corresponding operation is a tensor by the space $B(S)$, not an internal cartesian product unless the coefficient theory has that interpretation.

If $f:S'\to S$ is a weak homotopy equivalence, the natural homology comparison goes in the direction
\[
 p'_!p'^*a\longrightarrow p_!p^*a,
\]
and is represented here by $B(f)\times\mathrm{id}_a$. It is an isomorphism. The display on manuscript p.~132 prints the opposite direction. Once invertibility has been established one may of course write its inverse, but the map furnished directly by the functor $f$ is the one above. This directional slip does not invalidate the argument.

The connection with full faithfulness on constants is an application of adjunction and Yoneda. For coefficients $a,b$, one has
\[
 \mathrm{Hom}_{\HOT(S)}(p^*a,p^*b)
  \cong \mathrm{Hom}_{\mathrm{Hot}}(p_!p^*a,b).
\]
Restriction by $f$ becomes precomposition with the displayed homology comparison under this bijection. If that comparison is invertible for every $a$, restriction is fully faithful on constant objects. Conversely, bijectivity for every target coefficient $b$ makes the comparison invertible by Yoneda. Thus integration provides exactly the missing argument, rather than a merely suggestive analogy with ordinary homology.

This is the same general method already explained in DER\_I-a at manuscript pp.~61--66: express maps between constant objects by an adjunction and then use Yoneda. Its new content here is a concrete categorical model for the left adjoint. The derived Kan-extension results of \citet{Cisinski2003} put this method into a general model-category framework. They do not require collapsing coherent diagrams to diagrams in the ordinary homotopy category.
\end{dcomment}

\begin{dcomment}{exposition}{p137}{Why postcomposition already computes the direct image}
Grothendieck's proposed integration functor has a fortunate feature: it does not change the total category. For $f:T\to S$, let $L_f:\Cat/T\to\Cat/S$ send $q:X\to T$ to $fq:X\to S$. Coarse weak equivalences are preserved immediately. The substantive assertion is that $L_f$ also preserves $W_T^{\mathrm d}$, even when $f$ has no properness or smoothness hypothesis.

Fix $s\in S$ and form $U=s\backslash T$, with its forgetful map $r:U\to T$. This map is the base change of $s\backslash S\to S$, hence is a co-local isomorphism in the terminology of these pages. A $W_T^{\mathrm d}$-equivalence remains a coarse weak equivalence after such a base change. Its pullback is exactly the map
\[
 s\backslash X\longrightarrow s\backslash Y
\]
calculated using the composites to $S$. This holds for every $s$, proving the assertion for $W^{\mathrm d}$. Passing to opposite categories gives the $W^{\mathrm g}$ assertion. The proof uses stability of these elementary localizing maps, not stability of arbitrary pullback on all relative weak equivalences.

For completeness, the categorical stability used here is literal. Given a strict pullback square $A'=A\times_BB'$, and an object $a'=(a,b')$, an object of $a'\backslash A'$ is a pair of arrows $a\to x$ and $b'\to y'$ having equal image in $B$. Morphisms are likewise pairs of commuting triangles with equal image. Therefore
\[
 a'\backslash A'\cong
 (a\backslash A)\times_{b\backslash B}(b'\backslash B')
\]
as categories, on both objects and morphisms. Pulling back an isomorphism of the relevant undercategories gives another isomorphism. The single occurrence of ``isomorphisme local'' on p.~136 should be read as ``co-isomorphisme local'' in this undercategory calculation; the surrounding proof and the earlier paragraph determine the intended orientation.

Since $L_gL_f=L_{gf}$ before localization, the induced direct images compose canonically, and they commute with the further localization to coarse relative types. They do not generally preserve the locally constant subcategories. For an explicit example, take $j:e=\{0\}\to[1]$ and the point over $e$. Its direct image is $e\to[1]$ at $0$, representing $\varnothing\to e$, which is not locally constant. Its locally constant reflection is the constant point. Thus
\[
 \rho_S f_!\simeq f_!^{\mathrm{lc}}\rho_T
\]
holds, whereas the analogous square with the fully faithful inclusions need not commute.

The reservation in footnote~67 has a positive answer for $W_\infty$ when $f$ is a Hot-fibration. Indeed, for every $s\to s'$ in $S$, the transition $v:s'\backslash T\to s\backslash T$ is a weak equivalence. For a locally constant coefficient diagram $F$ on $T^{\mathrm{op}}$, pointwise left Kan extension expresses the corresponding transition in $f_!F$ as the homotopy-colimit comparison along $v^{\mathrm{op}}$, with the pulled-back locally constant coefficients. Such a comparison is invertible: restriction along a weak equivalence of indexing categories is an equivalence on locally constant diagrams, and the adjunction describing maps from a homotopy colimit into any constant coefficient, followed by Yoneda, proves invariance of the homotopy colimit. Thus every transition of $f_!F$ is invertible. This is exactly the presheaf-variance form of the formal Serre theorem \citep[Theorem~1.3 and Propositions~1.10, 1.15]{Cisinski2009LocallyConstant}; the comma-transition hypothesis explains why it applies to the Hot-fibrations of these pages.

One can already deduce the left adjoint on locally constant objects once the ordinary adjunction is established. If $i_T,i_S$ are the inclusions and $\rho_T,\rho_S$ their left adjoints, then
\[
 f_!^{\mathrm{lc}}\simeq\rho_S f_!i_T
\]
is left adjoint to $f^*_{\mathrm{lc}}$. Indeed, the adjunction bijections successively identify maps from the left side to $B$ with maps from $i_TA$ to $f^*i_SB\simeq i_Tf^*_{\mathrm{lc}}B$, and then with maps from $A$ to $f^*_{\mathrm{lc}}B$. This formal calculation explains why failure to preserve local constancy does not prevent a left adjoint inside the locally constant category. Grothendieck develops this point in the continuation, manuscript pp.~147--149; it does not repair the failure of the pointwise axiom demonstrated above.
\end{dcomment}

\begin{dcomment}{exposition}{p145}{The unit and counit, with both triangle identities}
Grothendieck now verifies that the geometric direct image is the promised left adjoint. The useful feature of his proof is that it exposes exactly where one may invert a replacement map. We give the verification in one notation, including the last step which falls just beyond the supplied transcription. Put $T=S'$, write $L=L_f$ for postcomposition, and write $R$ for derived restriction. Let $i_X:X\to\Phi_SX$ be the natural $W_S^{\mathrm d}$-equivalence. Then
\[
 R[X]=[\Phi_SX\times_ST].
\]
For a smooth representative $X$, one may identify this with $[X\times_ST]$, since both $X$ and $\Phi_SX$ are smooth and their comparison remains a relative weak equivalence after every base change.

For $q:Y\to T$, the unit is represented by the honest functor
\[
 a_Y=(i_{LY},q):Y\longrightarrow\Phi_S(LY)\times_ST.
\]
The pair lands in the fibre product because $i_{LY}$ is a map over $S$. Naturality of $i$ makes $a$ natural in $Y$. For an arbitrary $X\to S$, the counit is the localized morphism
\[
 b_X=[i_X]^{-1}[\pi_1]:[\Phi_SX\times_ST]\longrightarrow[X].
\]
Both factors are defined over $S$. Naturality follows from that of $i$ and $\pi_1$, and is retained on localization. If $X$ is smooth, this expression agrees with the projection $X\times_ST\to X$ under the comparison identifying the two models for $R[X]$.

The first triangle identity contains no hidden lifting argument. Since $\pi_1a_Y=i_{LY}$ as functors over $S$, one has
\[
 b_{LY}\,L(a_Y)=[i_{LY}]^{-1}[\pi_1a_Y]
  =[i_{LY}]^{-1}[i_{LY}]=\mathrm{id}_{LY}.
\]
This is exactly the corrected construction on manuscript p.~143. The abandoned attempt on p.~142 would have inverted $i_X\times_ST$. For an arbitrary $X$ this pullback need not be a weak equivalence; the example $\{1\}\to[1]$ pulled back to $0$ gives $\varnothing\to e$. Inverting $i_X$ over $S$ first avoids that error.

For the second triangle choose a fibred, hence smooth, representative $X\to S$ and put $Y=X\times_ST$. Write $h:Y\to X$ for the projection over $S$. The derived image $R(b_X)$ is represented by
\[
 [\Phi_S Y\times_ST]
 \xrightarrow{[\Phi_S(h)\times_ST]}
 [\Phi_S X\times_ST]
 \xrightarrow{[i_X\times_ST]^{-1}}
 [X\times_ST].
\]
The last inverse exists because $X$ and $\Phi_SX$ are smooth over $S$. To compute its composite with $a_Y$, compare the two projections of the maps into $\Phi_SX\times_ST$. Their first projections agree by naturality:
\[
 \Phi_S(h)i_Y=i_Xh.
\]
Their second projections are both the structural map $Y\to T$. The universal property of the fibre product therefore gives
\[
 (\Phi_S(h)\times_ST)a_Y=i_X\times_ST.
\]
Multiplying by the inverse just displayed proves $R(b_X)a_{RX}=\mathrm{id}_{RX}$. Since fibred representatives give an equivalence with the localized category, the identity holds for every object and, by naturality, as an identity of transformations. We have proved $f_!\dashv f^*$.

The last displayed triangle of the excerpt is precisely this equality. Grothendieck completes its proof on manuscript p.~146 by the same naturality square \citep[VI, ms.~pp.~146--147]{Grothendieck1990VI}; the end of the package is a cut in the proof, not a missing argument in the work. There is, however, a genuine typographical slip on p.~144: the target of $i_X\times_S\mathrm{id}_T$ must be $\Phi_S(X)\times_ST$, not $\Phi_S(X')\times_ST$. The following diagram in the edition already has the correct target. Likewise, both questioned curved arrows on pp.~141 and 144 are solid in the full edition; they express the composite being identified, not additional weak equivalences to be inverted.
\end{dcomment}
